\documentclass[12pt, a4paper]{article}
\usepackage[utf8]{inputenc}
\usepackage[T1]{fontenc}
\usepackage[ngermanb]{babel}
\usepackage{amsmath}
\usepackage{amssymb}
\usepackage{graphicx}
\usepackage{geometry}
\usepackage[hidelinks]{hyperref}
\usepackage{booktabs}
\usepackage{float}
\usepackage{tikz}
\usepackage{circuitikz}
\ctikzset{resistor=european}
\usepackage{pgfplots}
\usepackage{subcaption}
\usepackage[expansion=false]{microtype}
\usepackage{enumitem}
\usepackage{xcolor}
\usepackage{array}

% Farben fuer Halbleiter-Diagramme
\definecolor{ptype}{RGB}{255,180,100}
\definecolor{ntype}{RGB}{100,180,255}
\definecolor{depzone}{RGB}{220,220,220}
\definecolor{efeld}{RGB}{200,50,50}

% Manuelle Trennungshilfen
\hyphenation{Tech-no-lo-gie Halb-wel-len-gleich-rich-ter La-dungs-trä-ger-dich-te}

\setlist[itemize,1]{label=--}

\hypersetup{colorlinks=true,linkcolor=blue!60!black,urlcolor=blue!60!black,citecolor=blue!60!black}

\pgfplotsset{compat=1.18}

\usetikzlibrary{arrows.meta, positioning, calc, patterns,
	decorations.pathmorphing, decorations.markings, shadings, fadings}

\geometry{margin=2.5cm}

\title{\textbf{Physikalische Strukturen und\\ mathematische Beschreibbarkeit}\\
	\large Über die wechselseitige Ergänzung von Physik und Mathematik}
\author{Stephan Epp}
\date{8. März 2026}

\newcommand{\uititle}[1]{\textbf{$u$-$i$-Diagramm:} #1}

\begin{document}
	
	\maketitle
	\tableofcontents
	\newpage
	
	% ============================================================
	\section{Einleitung}
	% ============================================================
	
	Es gibt eine technische Anwendung für die Mathematik, in der das technische oder physikalische Verhalten \textit{physikalischer Strukturen} ideal durch die mathematische Beschreibung beschrieben wird. Mathematik und Physik ergänzen sich hier sehr gut und sehr anschaulich.
	
	Die zentrale Beobachtung dieser Arbeit ist, dass bestimmte physikalische Strukturen -- insbesondere elektronische Schaltungen und mechanische Systeme -- eine innere Konsistenz besitzen, die sich erst dann vollständig erschließt, wenn sowohl die physikalische als auch die mathematische Perspektive gemeinsam betrachtet werden. Eine rein mathematische Betrachtung kann dabei zu scheinbar widersprüchlichen oder physikalisch unsinnigen Aussagen führen, die sich erst durch das Hinzuziehen der physikalischen Realität auflösen. Umgekehrt vermag die Mathematik das physikalische Verhalten präzise zu erfassen und vorherzusagen -- sofern das zugrundeliegende physikalische Modell vollständig berücksichtigt wird.
	
	Wichtig in dieser Arbeit sind all jene physikalischen Strukturen, in denen sich Physik und Mathematik, wenn sie beide betrachtet werden, den vollständigen Sinn für die physikalische Struktur mit ihrer Beschreibung ergeben. Zu jedem betrachteten System werden in dieser erweiterten Fassung stets zwei Darstellungen gegeben: die elektronische Schaltung (gezeichnet mit der \texttt{circuitikz}-Bibliothek) sowie das zugehörige $u$-$i$-Diagramm, das den zeitlichen Verlauf von Spannung und Strom zeigt.
	
	Wenn man so will, löst sich die Spannung dieser ganz einfachen Struktur, bestehend aus einer Leitung und einem Widerstand, am Widerstand selber auf. Diese einfache physikalische Struktur ist die einfachste, die es gibt.

	\subsubsection*{Die einfachste physikalische Struktur: Leitung und Widerstand}

	Bevor die komplexeren Strukturen dieser Arbeit betrachtet werden, lohnt es sich, bei der allereinfachsten zu verweilen -- denn gerade in ihrer Schlichtheit offenbart sich das Grundprinzip am deutlichsten. Eine Leitung und ein Widerstand: mehr ist es nicht. Und doch enthält diese minimale Anordnung bereits den vollständigen Kern dessen, worum es in dieser Arbeit geht.

	Ein Widerstand $R$ ist ein passives Bauelement, das elektrische Energie in Wärme umwandelt. Er besitzt keine Fähigkeit, Energie zu speichern -- weder im elektrischen Feld wie ein Kondensator, noch im magnetischen Feld wie eine Spule. Er \textit{verbraucht} Energie, und zwar in dem Augenblick, in dem der Strom durch ihn hindurchfließt. Diese Eigenschaft -- die unmittelbare, verlustbehaftete Energieumwandlung -- macht ihn zur physikalisch klarsten und mathematisch einfachsten Struktur, die die Elektrotechnik kennt.

	\textbf{Das Ohmsche Gesetz als Kopplungsrelation.}
	Die physikalische Grundlage dieser Struktur ist das Ohmsche Gesetz:
	\begin{equation}
		U = R \cdot I,
		\label{eq:ohm}
	\end{equation}
	Das Ohmsche Gesetz ist zunächst nichts anderes als die Definition eines linearen Zusammenhangs zwischen Spannung $U$ und Strom $I$. Der Widerstand $R$ ist dabei die Proportionalitätskonstante. Diese Kopplung -- $U$ und $I$ sind nicht unabhängig, sondern durch $R$ miteinander verknüpft -- ist der entscheidende physikalische Inhalt.

	Würde man die Formel für die Leistung $P = U \cdot I$ rein mathematisch betrachten, ohne die physikalische Kopplung zu kennen, so könnte man denken: Für $U > 0$ und $I < 0$ ergibt sich $P < 0$ -- der Widerstand würde Energie erzeugen. Das ist mathematisch möglich, physikalisch aber ausgeschlossen. Denn das Ohmsche Gesetz erzwingt, dass $U$ und $I$ stets das gleiche Vorzeichen haben: Ein positiver Strom erzeugt eine positive Spannung, ein negativer Strom eine negative Spannung. Damit gilt stets:
	\begin{equation}
		P = U \cdot I = (R \cdot I) \cdot I = R \cdot I^2 \geq 0, \quad \forall\, I \in \mathbb{R},\; R \geq 0.
		\label{eq:leistung_widerstand_einl}
	\end{equation}
	Die quadratische Form $I^2$ ist der mathematische Ausdruck der physikalischen Tatsache, dass ein Widerstand keine Energie erzeugen kann. Die Nicht-Negativität der Leistung ist keine mathematische Zufälligkeit -- sie ist die direkte Konsequenz der physikalischen Kopplung zwischen Spannung und Strom.

	\textbf{Die Potenzialdifferenz als physikalische Ursache.}
	Betrachten wir die Struktur nun aus der Perspektive des elektrischen Potenzials. Ein elektrisches Potenzial $\varphi$ ist eine skalare Feldgröße, die jedem Punkt im Raum eine Energie pro Ladungseinheit zuweist. Die Spannung $U$ zwischen zwei Punkten ist nichts anderes als die Differenz der Potenziale:
	\begin{equation}
		U = \varphi_1 - \varphi_2.
		\label{eq:potential_diff}
	\end{equation}
	Diese Potenzialdifferenz ist die \textit{Ursache} des Stromflusses. Ohne sie -- also bei $\varphi_1 = \varphi_2$ -- fließt kein Strom. Mit ihr -- also bei $\varphi_1 \neq \varphi_2$ -- werden freie Elektronen im Leiter beschleunigt, bis sie auf Gitterbausteine treffen und dabei kinetische Energie in Wärme umwandeln. Genau diese Wechselwirkung zwischen Elektronen und Gitter ist der mikroskopische Ursprung des ohmschen Widerstands.

	Makroskopisch beschreibt das Ohmsche Gesetz diesen Mechanismus durch die einzige Größe $R$. Die Feinheiten der Mikrostruktur -- die Gitterkonstante, die Streuquerschnitte, die Fermi-Geschwindigkeit der Elektronen -- sind in $R$ zusammengefasst. Dies ist ein erstes, schönes Beispiel dafür, wie Mathematik (hier: eine lineare Gleichung) eine komplexe physikalische Wirklichkeit präzise und vollständig beschreibt.

	\textbf{Von links nach rechts: Richtung und Konvention.}
	Die Formulierung, dass die Potenzialdifferenz ``von links nach rechts betrachtet'' wird, ist mehr als eine räumliche Aussage -- sie betrifft die Frage der Konvention. In der Elektrotechnik wird der konventionelle Stromfluss als Bewegung positiver Ladungen definiert, also von hohem Potenzial ($\varphi_1$) zu niedrigem Potenzial ($\varphi_2$). Die tatsächliche Bewegung der Elektronen verläuft entgegengesetzt, von $\varphi_2$ nach $\varphi_1$ -- eine historische Konvention, die sich jedoch vollständig konsistent durch die gesamte Theorie zieht, sofern man sie konsequent anwendet.

	Die Wahl der Richtung ist mathematisch willkürlich, physikalisch aber bedeutsam: Sie legt fest, welches Vorzeichen der Strom hat, und damit, ob die Leistung $P = U \cdot I$ positiv oder negativ ist. Da $P \geq 0$ für einen ohmschen Widerstand gelten muss, ist das Vorzeichen des Stromes durch die Physik selbst festgelegt -- nicht durch Konvention.

	\textbf{Energiebilanz der einfachsten Struktur.}
	Die vollständige Energiebilanz dieser minimalen Struktur lässt sich wie folgt zusammenfassen: Eine Spannungsquelle (eine Batterie, ein Generator, ein beliebiger Energiewandler) stellt die Potenzialdifferenz $U = \varphi_1 - \varphi_2$ bereit. Diese treibt einen Strom $I = U/R$ durch den Widerstand. Pro Zeiteinheit wird dabei die Leistung
	\begin{equation}
		P = \frac{U^2}{R} = R \cdot I^2
	\end{equation}
	in Wärme umgewandelt. Die Quelle liefert genau diese Leistung -- sie gibt Energie ab, der Widerstand nimmt sie auf. Es gibt keine Speicherung, keine Rückkopplung, keine Dynamik. Die Struktur ist im stationären Gleichgewicht, sobald der Strom fließt.

	Diese Schlichtheit ist der Grund, warum der Widerstand die einfachste physikalische Struktur ist. Alle anderen Bauelemente -- Kondensator, Spule, Diode, Transistor -- fügen dieser einfachen Grundstruktur etwas hinzu: eine Zeitabhängigkeit, eine Nichtlinearität, eine Speicherfähigkeit. Der Widerstand allein kennt keine dieser Komplikationen. Er ist, in gewissem Sinne, die physikalische Verkörperung der linearen Algebra: Ein Eingang ($U$), ein Ausgang ($I$), eine Konstante ($R$), eine Gleichung ($U = R \cdot I$).

	\textbf{Warum diese Einfachheit der richtige Ausgangspunkt ist.}
	Es wäre ein Fehler, diese Einfachheit für Trivialität zu halten. Im Gegenteil: Gerade weil der Widerstand so einfach ist, kann man an ihm das Grundprinzip dieser Arbeit besonders klar zeigen. Physik und Mathematik stimmen hier vollständig überein -- nicht zufällig, sondern weil das Ohmsche Gesetz selbst \textit{die} Kopplungsrelation ist, die die mathematische Beschreibung mit der physikalischen Wirklichkeit verbindet. Alle anderen Strukturen dieser Arbeit sind, in einem tiefen Sinne, Variationen und Erweiterungen dieser einfachsten Kopplung. Sie fügen Energie hinzu, die gespeichert wird ($C$, $L$), sie fügen Nichtlinearitäten hinzu (Diode), sie fügen Steuerbarkeit hinzu (Transistor) -- aber das Grundprinzip bleibt: Die physikalische Kopplung zwischen den Größen erzwingt die mathematische Form der Beschreibung, und erst wenn beide zusammen betrachtet werden, erschließt sich der vollständige Sinn.

	\begin{figure}[h]
		\centering
		\begin{circuitikz}[thick]
			
			% Resistor starts at x=1.5, ends at x=4.5 (circuitikz R spans ~3 units)
			\draw (1,0)
			to[R, l=$R$, *-*] (3,0);
			
			% phi1 directly at left terminal of R
			\draw (1,0) node[below=4pt] {$\varphi_1$};
			
			% phi2 directly at right terminal of R
			\draw (3,0) node[below=4pt] {$\varphi_2$};
			
			% Arrow above with potential difference
			\draw[->, thick, blue] (1, 0.9) -- (3, 0.9)
			node[midway, above, blue] {$U = \varphi_1 - \varphi_2$};
		\end{circuitikz}
		\label{fig:resistor}
		\caption{Potentialdifferenz $U = \varphi_1 - \varphi_2$ am Widerstand $R$}
	\end{figure}
	Durch sie fließt ein Strom und es liegt eine Spannung an -- eine Potenzialdifferenz, die von links nach rechts betrachtet und beschrieben wird durch die beiden Unterschiede der Spannung $\varphi_1$ und $\varphi_2$. Die Abbildung zeigt dabei bereits das Wesentliche: Die Spannung $U = \varphi_1 - \varphi_2$ ist keine abstrakte Größe, sondern die physikalische Triebkraft, die den Strom durch den Widerstand erzwingt. Ohne diese Differenz kein Strom, ohne Strom keine Leistung, ohne Leistung keine Energieumwandlung -- die gesamte Physik dieser Struktur steckt in dieser einen skalaren Differenz zweier Potenzialwerte.

	Im Allgemeinen gilt für die Leistung:
	\begin{equation}
		P = R \cdot I^2
		\label{eq:leistung_R}
	\end{equation}
	Betrachte die physikalische Struktur oder die elektronische Schaltung zur Beschreibung der Leistung $P$. Im Allgemeinen ist $P$ abhängig sowohl von der Spannung als auch vom Strom, also
	\begin{equation}
		P = U \cdot I.
		\label{eq:leistung_UI}
	\end{equation}
	Würde man die Leistung für die Fälle $I > 0$ und $I < 0$ über der Zeit auftragen, würde sich für einen negativen Strom -- einmal theoretisch vorausgesetzt, dass $U$ in beiden Fällen positiv ist -- eine negative Leistung ergeben für den Fall $I < 0$.
	
	Betrachtet man nun aber weiter auch die physikalische Bedeutung dieser Aussage, muss ergänzt werden: Da $U$ und $I$ auch voneinander abhängig sind über den Widerstand $R$ -- nämlich durch das Ohmsche Gesetz $U = R \cdot I$ -- folgt, dass $P$ niemals negativ werden kann:
	\begin{equation}
		P = U \cdot I = (R \cdot I) \cdot I = R \cdot I^2 \geq 0,
		\quad \forall\, I \in \mathbb{R},\; R \geq 0.
	\end{equation}
	
	\begin{figure}[h]
		\centering
		\begin{subfigure}[b]{0.48\textwidth}
			\centering
			\begin{minipage}[c][6.5cm][c]{\linewidth}
				\centering
				\begin{circuitikz}[scale=1.0, transform shape]
					\draw (0,0) to[battery1, l=$U_0$, invert] (0,3)
					to[ammeter, l=$I$] (2,3)
					to[D, l=D] (4,3)
					to[R, l=$R$, v=$U_R$] (4,0)
					to[short] (0,0);
					\draw[->] (2.2,1.5) -- (2.8,1.5) node[midway,below]{$I \geq 0$};
				\end{circuitikz}
			\end{minipage}
			\caption{Schaltung: Diode sperrt $I < 0$}
		\end{subfigure}
		\hfill
		\begin{subfigure}[b]{0.48\textwidth}
			\centering
			\begin{minipage}[c][6.5cm][c]{\linewidth}
				\centering
				\begin{tikzpicture}
					\begin{axis}[
						% xlabel={Zeit $t$},
						ylabel={Amplitude},
						xmin=0, xmax=6.5,
						ymin=-1.5, ymax=1.5,
						width=7cm, height=5cm,
						grid=major,
						legend pos=north east,
						xtick={0,1.57,3.14,4.71,6.28},
						xticklabels={$0$,$\frac{T}{4}$,$\frac{T}{2}$,$\frac{3T}{4}$,$T$}
						]
						\addplot[blue, thick, domain=0:6.28, samples=100] {sin(deg(x))};
						\addlegendentry{$u(t) = U_0\sin(\omega t)$}
						\addplot[red, thick, domain=0:6.28, samples=200]
						{(sin(deg(x)) > 0) ? sin(deg(x)) : 0};
						\addlegendentry{$i(t) \geq 0$ (Diode)}
					\end{axis}
				\end{tikzpicture}
			\end{minipage}
			\caption{$u$-$i$-Diagramm: Diode sperrt negative Welle}
		\end{subfigure}
		\caption{Einfache Dioden-Widerstand-Schaltung und zugehöriges $u$-$i$-Diagramm. Die Diode erzwingt $I \geq 0$, sodass $P = R \cdot I^2 \geq 0$ stets gilt.}
		\label{fig:diode_schaltung}
	\end{figure}
	
	Die physikalische Einschränkung -- die Diode lässt nur eine Stromrichtung zu -- ist es, die die mathematische Bedingung $I \geq 0$ erzwingt und damit die Nicht-Negativität der Leistung auf Schaltungsebene sicherstellt. Das $u$-$i$-Diagramm zeigt deutlich: Die negative Halbwelle des Stromes existiert physikalisch nicht. Was die Mathematik für $I < 0$ ausrechnen würde, ist in der Realität durch die Diode blockiert.
	
	% ============================================================
	\section[Diode: Grundlage der CMOS-Technologie]{Diode: Grundlage der CMOS-Technologie}
	% ============================================================
	
	Die Diode ist die grundlegende physikalische Schaltung für die
	CMOS-Technologie. Die Diode ist ein Halbleiter.
	
	\textbf{Ein Halbleiter leitet nur in einem Fall, das heißt nur in der
		Hälfte aller Fälle. Zum Beispiel von links nach rechts, aber nicht von rechts
		nach links. Oder aber von rechts nach links, aber nicht von links nach rechts.}
	
	Diese fundamentale Aussage -- scheinbar einfach, in Wirklichkeit von tiefer
	physikalischer und mathematischer Bedeutung -- ist der Ausgangspunkt der
	gesamten modernen Halbleitertechnologie. Aus ihr folgt die Logik, die
	Rechenfähigkeit, und schließlich die gesamte digitale Welt. Im Folgenden wird
	diese Aussage physikalisch und mathematisch vollständig begründet.
	
	\subsection{Aufbau des Halbleiters: p-dotiertes und n-dotiertes Material}
	
	Ein reiner Halbleiter wie Silizium (Si) ist bei Raumtemperatur ein schlechter
	Leiter. Erst durch gezielte Verunreinigung -- die sogenannte \textit{Dotierung}
	-- wird er leitfähig. Dabei unterscheidet man zwei Typen:
	
	\textbf{n-dotiertes Silizium (Donatoren):} Durch Einbringen von
	Phosphor-Atomen (fünfwertig) in das Siliziumgitter entsteht ein Überschuss
	freier Elektronen. Diese negativen Ladungsträger sind die Majoriträtsladungsträger.
	Die positiv geladenen Atomrümpfe (Donatoren) sind ortsfest im Gitter verankert
	und können nicht fließen.
	
	\textbf{p-dotiertes Silizium (Akzeptoren):} Durch Einbringen von Bor-Atomen
	(dreiwertig) entsteht ein Mangel an Elektronen -- sogenannte
	\textit{Löcher} (Defektelektronen). Diese positiven Ladungsträger sind die
	Majoritätsladungsträger im p-Gebiet. Die negativ geladenen Akzeptor-Atomrümpfe
	sind ebenfalls ortsfest.
	
	\begin{figure}[h]
		\centering
		\begin{tikzpicture}[scale=1.0]
			% p-region background
			\fill[ptype!60] (0,0) rectangle (4,3);
			% n-region background
			\fill[ntype!60] (4,0) rectangle (8,3);
			% Labels
			\node[font=\Large\bfseries] at (2,2.5) {p-Gebiet};
			\node[font=\Large\bfseries] at (6,2.5) {n-Gebiet};
			% Holes in p-region (circles with +)
			\foreach \x in {0.5,1.2,1.9,2.6,3.3} {
				\foreach \y in {0.5,1.2,1.9} {
					\draw[thick, fill=white] (\x,\y) circle (0.18);
					\node[font=\scriptsize] at (\x,\y) {$+$};
				}
			}
			% Fixed negative ions in p-region
			\foreach \x in {0.85,1.55,2.25,2.95,3.65} {
				\foreach \y in {0.85,1.55} {
					\node[font=\scriptsize, color=blue!70!black] at (\x,\y) {\large$\ominus$};
				}
			}
			% Electrons in n-region (filled circles with -)
			\foreach \x in {4.7,5.4,6.1,6.8,7.5} {
				\foreach \y in {0.5,1.2,1.9} {
					\draw[thick, fill=gray!40] (\x,\y) circle (0.18);
					\node[font=\scriptsize, color=white] at (\x,\y) {$-$};
				}
			}
			% Fixed positive ions in n-region
			\foreach \x in {4.35,5.05,5.75,6.45,7.15} {
				\foreach \y in {0.85,1.55} {
					\node[font=\scriptsize, color=red!70!black] at (\x,\y) {\large$\oplus$};
				}
			}
			% Junction line
			\draw[very thick, dashed] (4,0) -- (4,3);
			\node[below] at (4,0.1) {pn-Ãœbergang};
			% Legend
			\draw[thick, fill=white] (0.3,-0.8) circle (0.15); \node[font=\scriptsize] at (0.3,-0.8) {$+$};
			\node[right, font=\small] at (0.55,-0.8) {Löcher (beweglich)};
			\draw[thick, fill=gray!40] (5.0,-0.8) circle (0.15); \node[font=\scriptsize, color=white] at (3.5,-0.8) {$-$};
			\node[right, font=\small] at (5.25,-0.8) {Elektronen (beweglich)};
			\node[font=\small, color=blue!70!black] at (0.3,-1.3) {\large$\ominus$};
			\node[right, font=\small] at (0.55,-1.3) {Akzeptor-Ion (ortsfest)};
			\node[font=\small, color=red!70!black] at (5.0,-1.3) {\large$\oplus$};
			\node[right, font=\small] at (5.25,-1.3) {Donator-Ion (ortsfest)};
			% border
			\draw[thick] (0,0) rectangle (8,3);
		\end{tikzpicture}
		\caption{Schematischer Querschnitt durch ein p-dotiertes (links) und n-dotiertes
			(rechts) Halbleitermaterial vor dem Zusammenfügen. Die beweglichen
			Ladungsträger (Löcher und Elektronen) sowie die ortsfesten Gitterionen sind
			dargestellt.}
		\label{fig:pn_getrennt}
	\end{figure}
	
	\subsection{Die Raumladungszone: Entstehung und physikalische Bedeutung}
	
	Wenn p-dotiertes und n-dotiertes Material direkt miteinander in Kontakt gebracht
	werden, entsteht ein pn-Ãœbergang -- die grundlegende Struktur der Diode. An
	der Grenzfläche zwischen p- und n-Gebiet beginnt ein Diffusionsprozess: Die
	freien Elektronen aus dem n-Gebiet diffundieren in das p-Gebiet, und die Löcher
	aus dem p-Gebiet diffundieren in das n-Gebiet. Dieser Vorgang gehorcht dem
	Fickschen Diffusionsgesetz:
	\begin{equation}
		J_{\mathrm{diff}} = -D \cdot \frac{dn}{dx},
	\end{equation}
	wobei $D$ die Diffusionskonstante und $dn/dx$ der Konzentrationsgradient der
	Ladungsträger ist.
	
	Da die Elektronen aus dem n-Gebiet in das p-Gebiet diffundieren, hinterlassen
	sie positiv geladene Donator-Ionen im n-Gebiet. Ebenso hinterlassen die ins
	n-Gebiet diffundierenden Löcher negativ geladene Akzeptor-Ionen im p-Gebiet.
	In der unmittelbaren Umgebung des Übergangs sind alle freien Ladungsträger
	rekombiniert -- es verbleiben nur die ortsfesten Ionen. Dieser Bereich wird als
	\textbf{Raumladungszone} (auch: Sperrschicht, Verarmungszone, engl.:
	\textit{depletion region}) bezeichnet.
	
	\begin{figure}[h]
		\centering
		\begin{tikzpicture}[scale=1.0]
			% ---- Zones ----
			% p-neutral zone
			\fill[ptype!50] (0,0.5) rectangle (2.5,4.5);
			% depletion in p-side
			\fill[depzone] (2.5,0.5) rectangle (4.0,4.5);
			% depletion in n-side
			\fill[depzone] (4.0,0.5) rectangle (5.5,4.5);
			% n-neutral zone
			\fill[ntype!50] (5.5,0.5) rectangle (8,4.5);
			
			% Zone labels
			\node[font=\bfseries] at (1.25,4.9) {p-neutral};
			\node[font=\bfseries] at (3.25,4.9) {};
			\node[font=\bfseries] at (4.75,4.9) {};
			\node[font=\bfseries] at (6.75,4.9) {n-neutral};
			\node[font=\small] at (4.0,5.2) {\textbf{}};
			\draw[<->, thick] (2.5,5.0) -- (5.5,5.0);
			
			% Holes in p-neutral region
			\foreach \x in {0.4,0.9,1.4,1.9} {
				\foreach \y in {1.0,1.8,2.6,3.4} {
					\draw[fill=white, thick] (\x,\y) circle (0.15);
					\node[font=\tiny] at (\x,\y) {$+$};
				}
			}
			% Negative acceptor ions in p-depletion (RLZ)
			\foreach \x in {2.7,3.1,3.5,3.9} {
				\foreach \y in {1.0,1.8,2.6,3.4} {
					\node[font=\normalsize, color=blue!80!black] at (\x,\y) {$\ominus$};
				}
			}
			% Positive donor ions in n-depletion (RLZ)
			\foreach \x in {4.1,4.5,4.9,5.3} {
				\foreach \y in {1.0,1.8,2.6,3.4} {
					\node[font=\normalsize, color=red!80!black] at (\x,\y) {$\oplus$};
				}
			}
			% Electrons in n-neutral region
			\foreach \x in {5.7,6.2,6.7,7.2,7.7} {
				\foreach \y in {1.0,1.8,2.6,3.4} {
					\draw[fill=gray!50, thick] (\x,\y) circle (0.15);
					\node[font=\tiny, color=white] at (\x,\y) {$-$};
				}
			}
			
			% Electric field arrows (pointing from + to -, i.e. n->p inside RLZ)
			\foreach \y in {1.4, 2.2, 3.0} {
				\draw[->, very thick, color=efeld] (5.2,\y) -- (2.8,\y);
			}
			\node[color=efeld, font=\bfseries] at (4.0,-0.2) {$\vec{E}$ (eingebautes Feld)};
			
			% Diffusion arrows
			\draw[->, thick, blue!70!black, dashed] (5.5,4.3) -- (2.5,4.3)
			node[midway, above, font=\small] {Elektronen diffundieren};
			\draw[->, thick, orange!80!black, dashed] (2.5,0.7) -- (5.5,0.7)
			node[midway, below, font=\small] {Löcher diffundieren};
			
			% Junction line
			\draw[very thick] (4.0,0.5) -- (4.0,4.5);
			
			% Border
			\draw[thick] (0,0.5) rectangle (8,4.5);
			\node[below] at (4.0,0.4) {};
		\end{tikzpicture}
		\caption{Querschnitt durch den pn-Ãœbergang mit ausgebildeter Raumladungszone
			(RLZ). In der RLZ befinden sich keine freien Ladungsträger, nur die ortsfesten
			Gitterionen. Das eingebaute elektrische Feld $\vec{E}$ zeigt vom n-Gebiet zum
			p-Gebiet und wirkt der weiteren Diffusion entgegen.}
		\label{fig:raumladungszone}
	\end{figure}
	
	Die Raumladungszone ist elektrisch geladen: Die p-seitige Hälfte ist negativ
	(Akzeptoren), die n-seitige Hälfte ist positiv (Donatoren). Dadurch entsteht
	ein eingebautes elektrisches Feld $\vec{E}$, das vom n-Gebiet zum p-Gebiet
	zeigt -- also der Diffusion entgegen. Im Gleichgewichtszustand (ohne äußere
	Spannung) heben sich Diffusionsstrom und Driftstrom genau auf:
	\begin{equation}
		J_{\mathrm{diff}} + J_{\mathrm{drift}} = 0 \quad \Rightarrow \quad J_{\mathrm{ges}} = 0.
	\end{equation}
	
	Die zugehörige eingebaute Spannung (Diffusionsspannung) beträgt für Silizium
	typischerweise $U_D \approx 0{,}6\,\mathrm{V}$:
	\begin{equation}
		U_D = \frac{k_B T}{q} \ln\!\left(\frac{N_A \cdot N_D}{n_i^2}\right),
	\end{equation}
	wobei $k_B$ die Boltzmann-Konstante, $T$ die absolute Temperatur, $q$ die
	Elementarladung, $N_A$ und $N_D$ die Akzeptor- und Donatorkonzentrationen und
	$n_i$ die intrinsische Ladungsträgerdichte sind.
	
	\subsection{Die Diode in Durchlassrichtung (Vorwärtspolung)}
	
	Wird eine positive Spannung $U_F > 0$ an die Diode angelegt (Plus an p, Minus
	an n), so wird das eingebaute elektrische Feld geschwächt. Die Raumladungszone
	wird schmaler. Ab einer Schwellspannung von $U_{th} \approx 0{,}6\,\mathrm{V}$
	(Silizium) ist die Raumladungszone so weit abgebaut, dass ein signifikanter
	Strom fließen kann.
	
	\begin{figure}[h]
		\centering
		\begin{tikzpicture}[scale=0.95]
			% p-neutral zone
			\fill[ptype!50] (0,0.5) rectangle (3.0,4.0);
			% depletion narrow
			\fill[depzone] (3.0,0.5) rectangle (3.8,4.0);
			\fill[depzone] (3.8,0.5) rectangle (4.6,4.0);
			% n-neutral zone
			\fill[ntype!50] (4.6,0.5) rectangle (7.5,4.0);
			
			% RLZ label (narrow)
			% \node[font=\small\bfseries] at (3.8,4.3) {};
			% \draw[<->, thick] (3.0,4.15) -- (4.6,4.15);
			
			% Holes in p
			\foreach \x in {0.5,1.1,1.7,2.3,2.8} {
				\foreach \y in {1.0,1.8,2.6,3.4} {
					\draw[fill=white, thick] (\x,\y) circle (0.15);
					\node[font=\tiny] at (\x,\y) {$+$};
				}
			}
			% Ions in RLZ
			\foreach \x in {3.2,3.6} {
				\foreach \y in {1.0,1.8,2.6,3.4} {
					\node[font=\normalsize, color=blue!80!black] at (\x,\y) {$\ominus$};
				}
			}
			\foreach \x in {4.0,4.4} {
				\foreach \y in {1.0,1.8,2.6,3.4} {
					\node[font=\normalsize, color=red!80!black] at (\x,\y) {$\oplus$};
				}
			}
			% Electrons in n
			\foreach \x in {4.9,5.4,5.9,6.4,6.9} {
				\foreach \y in {1.0,1.8,2.6,3.4} {
					\draw[fill=gray!50, thick] (\x,\y) circle (0.15);
					\node[font=\tiny, color=white] at (\x,\y) {$-$};
				}
			}
			
			% External voltage source
			\draw[thick] (-1.0,0.5) -- (-1.0,4.0);
			\draw[thick] (-1.0,4.0) -- (0,4.0) -- (0,3.8);
			\draw[thick] (-1.0,0.5) -- (0,0.5) -- (0,0.7);
			\node at (-1.0,2.25) [left, font=\small] {$+$};
			\node at (-1.0,2.25) [right, font=\small] {$U_F$};
			% Voltage label
			\draw[thick] (7.5,4.0) -- (8.5,4.0) -- (8.5,0.5) -- (7.5,0.5);
			\node at (8.5,2.25) [right, font=\small] {$-$};
			
			% Current arrows (conventional current p->n externally, n->p internally)
			\draw[->, very thick, red!80!black] (0.3,4.2) -- (7.2,4.2)
			node[midway, above, font=\bfseries\small] {$I_F \to$ (Durchlassstrom)};
			
			% Electric field (weakened, short arrows)
			\draw[->, thick, efeld] (4.3,2.2) -- (3.3,2.2);
			\node[color=efeld, font=\small] at (3.8,1.6) {$\vec{E}$ (geschwächt)};
			
			\draw[thick] (0,0.5) rectangle (7.5,4.0);
			\draw[very thick] (3.8,0.5) -- (3.8,4.0);
			\node[below, font=\small] at (3.8,0.55) {pn-Ãœbergang};
		\end{tikzpicture}
		\caption{Diode in Durchlassrichtung: Die externe Spannung $U_F$ schwächt das
			eingebaute Feld, die Raumladungszone wird schmaler, ein Durchlassstrom $I_F$
			fließt.}
		\label{fig:diode_forward}
	\end{figure}
	
	\subsection{Die Diode in Sperrrichtung (Rückwärtspolung)}
	
	Wird eine negative Spannung $U_R < 0$ angelegt (Plus an n, Minus an p),
	so wird das eingebaute Feld verstärkt. Die Raumladungszone wächst. Es fließt
	nur ein vernachlässigbar kleiner Sättigungssperrstrom $I_S$ (typisch
	$\approx \mathrm{nA}$ bis $\mu\mathrm{A}$).
	
	\begin{figure}[h]
		\centering
		\begin{tikzpicture}[scale=0.95]
			% p-neutral zone (smaller)
			\fill[ptype!50] (0,0.5) rectangle (1.5,4.0);
			% depletion (wide)
			\fill[depzone] (1.5,0.5) rectangle (3.8,4.0);
			\fill[depzone] (3.8,0.5) rectangle (6.0,4.0);
			% n-neutral zone (smaller)
			\fill[ntype!50] (6.0,0.5) rectangle (7.5,4.0);
			
			% RLZ label (wide)
			\node[font=\small\bfseries] at (3.75,4.3) {};
			\draw[<->, thick] (1.5,4.15) -- (6.0,4.15);
			
			% Holes in p
			\foreach \x in {0.4,0.9,1.4} {
				\foreach \y in {1.0,1.8,2.6,3.4} {
					\draw[fill=white, thick] (\x,\y) circle (0.15);
					\node[font=\tiny] at (\x,\y) {$+$};
				}
			}
			% Neg ions in p-RLZ
			\foreach \x in {1.7,2.1,2.5,2.9,3.3,3.7} {
				\foreach \y in {1.0,1.8,2.6,3.4} {
					\node[font=\normalsize, color=blue!80!black] at (\x,\y) {$\ominus$};
				}
			}
			% Pos ions in n-RLZ
			\foreach \x in {3.9,4.3,4.7,5.1,5.5,5.9} {
				\foreach \y in {1.0,1.8,2.6,3.4} {
					\node[font=\normalsize, color=red!80!black] at (\x,\y) {$\oplus$};
				}
			}
			% Electrons in n
			\foreach \x in {6.2,6.7,7.2} {
				\foreach \y in {1.0,1.8,2.6,3.4} {
					\draw[fill=gray!50, thick] (\x,\y) circle (0.15);
					\node[font=\tiny, color=white] at (\x,\y) {$-$};
				}
			}
			% Strong E-field arrows
			\foreach \y in {1.2, 2.0, 2.8, 3.6} {
				\draw[->, very thick, efeld] (5.7,\y) -- (1.8,\y);
			}
			\node[color=efeld, font=\bfseries\small] at (3.75,0.1) {$\vec{E}$ (verstärkt)};
			
			% External reverse voltage
			\draw[thick] (-0.8,0.5) -- (-0.8,4.0);
			\draw[thick] (-0.8,4.0) -- (0,4.0);
			\draw[thick] (-0.8,0.5) -- (0,0.5);
			\node at (-0.8,2.25) [left, font=\small] {$-$};
			\node[font=\small] at (-0.3,2.25) {$U_R$};
			\draw[thick] (7.5,4.0) -- (8.3,4.0) -- (8.3,0.5) -- (7.5,0.5);
			\node at (8.3,2.25) [right, font=\small] {$+$};
			
			% Tiny reverse current
			\draw[->, thick, gray] (7.3,4.2) -- (0.2,4.2)
			node[midway, above, font=\small] {$I_S \approx 0$ (Sperrstrom)};
			
			\draw[thick] (0,0.5) rectangle (7.5,4.0);
			\draw[very thick] (3.8,0.5) -- (3.8,4.0);
			\node[below, font=\small] at (3.8,0.4) {};
		\end{tikzpicture}
		\caption{Diode in Sperrrichtung: Die externe Spannung verstärkt das eingebaute
			Feld, die Raumladungszone wächst stark an. Es fließt praktisch kein Strom.
			Dies ist die physikalische Ursache der Einwegigkeit des Halbleiters.}
		\label{fig:diode_reverse}
	\end{figure}
	
	\subsection{Die Shockley-Gleichung: Mathematische Beschreibung der Diodenkennlinie}
	
	Die vollständige mathematische Beschreibung des Diodenstroms liefert die
	\textbf{Shockley-Gleichung} (ideale Diodenkennlinie):
	\begin{equation}
		I_D(U) = I_S \left( e^{U / (n \cdot U_T)} - 1 \right),
		\label{eq:shockley}
	\end{equation}
	wobei
	\begin{itemize}
		\item $I_S$ der Sättigungssperrstrom (typisch $10^{-12}$ bis $10^{-6}\,\mathrm{A}$),
		\item $n$ der Idealitätsfaktor ($1 \leq n \leq 2$ für reale Dioden),
		\item $U_T = k_B T / q \approx 26\,\mathrm{mV}$ bei Raumtemperatur die thermische Spannung
	\end{itemize}
	sind.
	
	\textbf{Analyse der Shockley-Gleichung:}
	
	\textit{Für $U \gg U_T$ (Durchlassrichtung):}
	\begin{equation}
		I_D \approx I_S \cdot e^{U/(nU_T)} \gg 0.
	\end{equation}
	Der Strom steigt exponentiell an. Bereits eine kleine positive Spannung
	erzeugt einen großen Strom.
	
	\textit{Für $U \ll 0$ (Sperrrichtung):}
	\begin{equation}
		I_D \approx I_S \left(0 - 1\right) = -I_S.
	\end{equation}
	Der Strom sättigt bei dem sehr kleinen negativen Wert $-I_S$. Physikalisch
	bedeutet dies: In Sperrrichtung fließt praktisch kein Strom.
	
	\begin{figure}[h]
		\centering
		\begin{subfigure}[b]{0.48\textwidth}
			\centering
			\begin{minipage}[c][7.0cm][c]{\linewidth}
				\centering
				\begin{circuitikz}[scale=1.1, transform shape]
					\draw (0,0) to[battery1, invert, l=$U_F$] (0,2.5)
					to[short] (1,2.5)
					to[D] (3,2.5)
					to[R, l=$R$, v=$U_R$] (3,0)
					to[short] (0,0);
					% \draw[->, thick, red] (1.0,2.7) -- (2.0,2.7) node[midway,above]{$I_D$};
					\node[below] at (1.5,-0.3) {Durchlassschaltung};
				\end{circuitikz}
			\end{minipage}
			\caption{Schaltung: Diode in Durchlassrichtung}
		\end{subfigure}
		\hfill
		\begin{subfigure}[b]{0.48\textwidth}
			\centering
			\begin{minipage}[c][7.0cm][c]{\linewidth}
				\centering
				\begin{tikzpicture}
					\begin{axis}[
						% xlabel={Diodenspannung $U_D$ [V]},
						ylabel={Diodenstrom $I_D$ [mA]},
						xmin=-1.0, xmax=0.8,
						ymin=-0.2, ymax=5,
						width=7.5cm, height=5.5cm,
						grid=major,
						legend pos=north west,
						]
						% Shockley: I = Is*(exp(U/Ut)-1), Is=1e-12, Ut=0.026, scale to mA
						% For plot: I[mA] = 1e-9 * (exp(U/0.026)-1)
						\addplot[blue, very thick, domain=-1.0:0.55, samples=30] {0.001};
						\addplot[blue, very thick, domain=0.55:0.72, samples=60]
						{min(5.0, 0.001*exp((x-0.55)/0.026))};
						\addlegendentry{$I_D(U_D)$ Shockley}
						\addplot[dashed, gray, domain=-1.0:0.0] {0};
						\addplot[only marks, mark=*, red] coordinates {(0.6, 0.0)};
						\node[right, font=\small] at (axis cs: 0.62, 0.3) {$U_{th}\approx 0{,}6\,\text{V}$};
						\draw[<->, thick, orange] (axis cs:-0.05, -0.18) -- (axis cs:-1.0,-0.18)
						node[midway, below, font=\tiny] {$-I_S \approx 0$};
					\end{axis}
				\end{tikzpicture}
			\end{minipage}
			\caption{Diodenkennlinie nach Shockley}
		\end{subfigure}
		\caption{Diodenschaltung und Kennlinie: Die exponentielle Shockley-Gleichung
			beschreibt das stark asymmetrische Verhalten. In Durchlassrichtung ($U > 0{,}6\,\mathrm{V}$)
			fließt exponentiell wachsender Strom, in Sperrrichtung ($U < 0$) praktisch
			keiner.}
		\label{fig:diode_kennlinie}
	\end{figure}
	
	\subsection{Breite der Raumladungszone: mathematische Beschreibung}
	
	Die Breite der Raumladungszone $d$ ist abhängig von der angelegten Spannung:
	\begin{equation}
		d(U) = \sqrt{\frac{2\varepsilon_r \varepsilon_0}{q} \left(\frac{1}{N_A} +
			\frac{1}{N_D}\right)(U_D - U)},
		\label{eq:rlz_breite}
	\end{equation}
	wobei $\varepsilon_r \approx 11{,}7$ für Silizium, $U_D$ die eingebaute
	Diffusionsspannung und $U$ die externe Spannung ist.
	
	Für $U > 0$ (Durchlassrichtung): $d$ wird kleiner $\to$ Barriere sinkt $\to$
	Strom steigt.
	
	Für $U < 0$ (Sperrrichtung): $d$ wird größer $\to$ Barriere wächst $\to$
	kein Strom.
	
	\begin{figure}[h]
		\centering
		\begin{tikzpicture}
			\begin{axis}[
				xlabel={Externe Spannung $U$ [V]},
				ylabel={RLZ-Breite $d$ [normiert]},
				xmin=-3, xmax=0.6,
				ymin=0, ymax=3.5,
				width=10cm, height=5cm,
				grid=major,
				legend pos=north east,
				]
				% d ~ sqrt(U_D - U), U_D = 0.7
				\addplot[blue, very thick, domain=-3:0.65, samples=200]
				{sqrt(0.7 - x)};
				\addlegendentry{$d \propto \sqrt{U_D - U}$}
				\addplot[only marks, mark=*, red] coordinates {(0,0.837)};
				\node[right, font=\small] at (axis cs:0.05,0.95) {$d_0$ (Gleichgewicht)};
				\addplot[dashed, orange, domain=-3:0.6] {0.837};
				\node[left, font=\small, color=orange] at (axis cs:-3,0.95) {$d_0$};
			\end{axis}
		\end{tikzpicture}
		\caption{Breite der Raumladungszone als Funktion der angelegten Spannung. In
			Sperrrichtung ($U < 0$) wächst $d$ proportional zu $\sqrt{U_D - U}$, in
			Durchlassrichtung ($U \to U_D$) geht $d \to 0$.}
		\label{fig:rlz_breite}
	\end{figure}
	
	\subsection{Potentialverlauf und Banddiagramm des pn-Ãœbergangs}
	
	Der Potentialverlauf über den pn-Übergang ist eine direkte mathematische
	Konsequenz der Raumladung. Die Poisson-Gleichung lautet:
	\begin{equation}
		\frac{d^2 \varphi}{dx^2} = -\frac{\rho(x)}{\varepsilon_r \varepsilon_0},
	\end{equation}
	wobei $\rho(x)$ die Raumladungsdichte ist:
	\begin{equation}
		\rho(x) = \begin{cases}
			-q N_A & \text{für } -d_p \leq x < 0 \quad \text{(p-seitige RLZ)}\\
			+q N_D & \text{für } 0 < x \leq d_n \quad \text{(n-seitige RLZ)}\\
			0 & \text{sonst (neutrales Gebiet)}
		\end{cases}
	\end{equation}
	
	\begin{figure}[h]
		\centering
		\begin{tikzpicture}[scale=1.0]
			% Axes
			\begin{scope}
				\begin{axis}[
					name=rho,
					xlabel={},
					ylabel={Raumladung $\rho(x)$},
					xmin=-3, xmax=3,
					ymin=-1.5, ymax=1.5,
					width=12cm, height=3.5cm,
					xtick={-2,0,2}, xticklabels={$-d_p$,$0$,$d_n$},
					ytick={-1,0,1},
					yticklabels={$-qN_A$,$0$,$+qN_D$},
					grid=none,
					axis lines=center,
					title={Raumladungsdichte},
					]
					\addplot[blue, very thick] coordinates {(-3,0)(-2,0)(-2,-1)(-0.01,-1)(-0.01,0)};
					\addplot[red, very thick] coordinates {(0.01,0)(0.01,1)(2,1)(2,0)(3,0)};
					\addplot[gray, dashed] coordinates {(-2,-1.4)(-2,1.4)};
					\addplot[gray, dashed] coordinates {(2,-1.4)(2,1.4)};
				\end{axis}
			\end{scope}
			
			\begin{scope}[yshift=-4.2cm]
				\begin{axis}[
					name=phi,
					xlabel={Ortskoordinate $x$},
					ylabel={Potential $\varphi(x)$},
					xmin=-3, xmax=3,
					ymin=-0.2, ymax=1.2,
					width=12cm, height=3.5cm,
					xtick={-2,0,2}, xticklabels={$-d_p$,$0$,$d_n$},
					grid=none,
					axis lines=center,
					title={Elektrisches Potential},
					]
					% parabola in p-depl, parabola in n-depl, flat outside
					% phi_p-side: x from -2 to 0: phi = 0.35*(x+2)^2/4
					\addplot[blue!70!black, very thick, domain=-3:-2, samples=5] {0};
					\addplot[blue!70!black, very thick, domain=-2:0, samples=50]
					{0.35*(x+2)*(x+2)/4};
					\addplot[red!70!black, very thick, domain=0:2, samples=50]
					{0.35 + 0.65*(1-(x-2)*(x-2)/4)};
					\addplot[red!70!black, very thick, domain=2:3, samples=5] {1.0};
					\addplot[gray, dashed] coordinates {(-2,-0.2)(-2,1.2)};
					\addplot[gray, dashed] coordinates {(2,-0.2)(2,1.2)};
					\node[right, font=\small] at (axis cs:2.1,1.05) {$U_D$};
				\end{axis}
			\end{scope}
		\end{tikzpicture}
		\caption{Raumladungsdichte $\rho(x)$ (oben) und resultierender Potentialverlauf
			$\varphi(x)$ (unten) über den pn-Übergang. Die Potentialbarriere $U_D$ muss
			überwunden werden, damit Ladungsträger den Übergang passieren können.}
		\label{fig:banddiagramm}
	\end{figure}
	
	\subsection{Die Diode als Grundlage der CMOS-Technologie}
	
	Die CMOS-Technologie (Complementary Metal-Oxide-Semiconductor) basiert auf
	zwei komplementären Transistortypen: dem NMOS- und dem PMOS-Transistor. Beide
	sind im Kern nichts anderes als gesteuerte pn-übergänge -- also gesteuerte
	Dioden. Die Steuerung erfolgt über eine Gate-Elektrode, die über eine dünne
	Oxidschicht ($\mathrm{SiO_2}$, typisch $< 10\,\mathrm{nm}$ in modernen Prozessen)
	vom Halbleiter isoliert ist.
	
	\begin{figure}[h]
		\centering
		\begin{tikzpicture}[scale=1.0]
			% Substrate p-type
			\fill[ptype!40] (0,0) rectangle (8,1.5);
			\node[font=\small\bfseries] at (4,0.3) {p-Substrat};
			
			% n+ source and drain
			\fill[ntype!80] (0.5,0.8) rectangle (2.0,1.5);
			\node[font=\scriptsize \bfseries] at (1.25,1.1) {n$^+$Source};
			\fill[ntype!80] (6.0,0.8) rectangle (7.5,1.5);
			\node[font=\scriptsize \bfseries] at (6.75,1.1) {n$^+$Drain};
			
			% Channel region (depletion or inversion)
			\fill[depzone!80] (2.0,0.8) rectangle (6.0,1.5);
			\node[font=\small] at (4.0,1.1) {Kanal (gesteuert)};
			
			% Gate oxide
			\fill[yellow!60] (2.0,1.5) rectangle (6.0,1.8);
			\node[font=\small] at (4.0,1.65) {};
			
			% Gate electrode
			\fill[gray!60] (2.0,1.8) rectangle (6.0,2.2);
			\node[font=\small\bfseries] at (4.0,2.0) {Gate};
			
			% Source contact
			\fill[gray!70] (0.5,1.5) rectangle (1.5,2.0);
			\node[above, font=\small] at (1.0,2.0) {S};
			% Drain contact
			\fill[gray!70] (6.5,1.5) rectangle (7.5,2.0);
			\node[above, font=\small] at (7.0,2.0) {D};
			
			% Gate voltage arrow
			\draw[->, thick] (4.0,2.6) -- (4.0,2.3) node[above, at start, font=\small]{$U_{GS}$};
			
			% Current arrow in channel
			\draw[->, very thick, red!80!black] (2.2,1.2) -- (5.8,1.2)
			node[midway, above, font=\small]{$I_{DS}$ (wenn $U_{GS} > U_{th}$)};
			
			% Diode symbols at junctions (Source-Substrat, Drain-Substrat)
			\draw[thick] (2.0,0) -- (2.0,0.8);
			\node[font=\tiny, rotate=90] at (2.0,0.4) {pn-Diode};
			\draw[thick] (6.0,0) -- (6.0,0.8);
			\node[font=\tiny, rotate=90] at (6.0,0.4) {pn-Diode};
			
			\draw[thick] (0,0) rectangle (8,2.2);
			\node[below] at (4.0,-0.3) {NMOS-Transistor: Querschnitt};
		\end{tikzpicture}
		\caption{Querschnitt eines NMOS-Transistors. Source und Drain sind
			n$^+$-dotierte Gebiete im p-Substrat -- sie bilden je eine pn-Diode mit dem
			Substrat. Das Gate steuert über das Oxid den Kanal zwischen Source und Drain.
			Ohne ausreichende Gate-Spannung ($U_{GS} < U_{th}$) sperren die pn-Dioden --
			genau wie eine einzelne Diode in Sperrrichtung.}
		\label{fig:nmos_querschnitt}
	\end{figure}
	
	Der NMOS-Transistor kann als \textbf{spannungsgesteuerter Schalter} betrachtet
	werden: Für $U_{GS} < U_{th}$ sperrt er (kein Kanal, pn-Diode in Sperrrichtung),
	für $U_{GS} > U_{th}$ leitet er (invertierter Kanal). Die pn-Diode zwischen
	Drain/Source und Substrat ist dabei stets in Sperrrichtung gepolt -- genau
	jenes Verhalten, das im vorangegangenen Abschnitt beschrieben wurde:
	
	Die Diode ist die physikalische Grundlage, auf der die gesamte Logik
	digitaler Schaltungen aufbaut. Ein CMOS-Inverter, ein NAND-Gatter, ein
	Flip-Flop, ein Mikroprozessor -- alle diese Strukturen sind letztlich
	Anordnungen von Transistoren, deren Kern der pn-Ãœbergang ist.
	
	% ============================================================
	\section{Der elektrische Kondensator}
	% ============================================================
	
	\subsection{Physikalische Beschreibung und Schaltung}
	
	Die im Kondensator gespeicherte elektrische Energie beträgt:
	\begin{equation}
		W_C = \frac{1}{2} C \cdot U^2 \geq 0
		\label{eq:kondensator_energie}
	\end{equation}
	
	Der Ladestrom beim Laden eines Kondensators über einen Widerstand lautet:
	\begin{equation}
		i_C(t) = \frac{U_0}{R} \cdot e^{-t/(RC)}, \qquad
		u_C(t) = U_0 \left(1 - e^{-t/(RC)}\right)
	\end{equation}
	
	\begin{figure}[h]
		\centering
		\begin{subfigure}[b]{0.48\textwidth}
			\centering
			\begin{minipage}[c][6.5cm][c]{\linewidth}
				\centering
				\begin{circuitikz}[scale=1.0, transform shape]
					\draw (0,0) to[battery1, l=$U_0$, invert] (0,3)
					to[R, l=$R$] (3,3)
					to[C, l=$C$, v=$u_C$] (3,0)
					to[short] (0,0);
					\node[draw, circle, inner sep=1pt] at (1.5, 1.5) {};
					% \draw[->, thick] (1.5,2.8) -- (1.5,2.2) node[right]{$i_C$};
				\end{circuitikz}
			\end{minipage}
			\caption{RC-Ladeschaltung}
		\end{subfigure}
		\hfill
		\begin{subfigure}[b]{0.48\textwidth}
			\centering
			\begin{minipage}[c][6.5cm][c]{\linewidth}
				\centering
				\begin{tikzpicture}
					\begin{axis}[
						% xlabel={Zeit $t / \tau$},
						ylabel={Normierte Größe},
						xmin=0, xmax=5,
						ymin=0, ymax=1.1,
						width=7cm, height=5cm,
						grid=major,
						legend pos=north east,
						]
						\addplot[blue, thick, domain=0:5, samples=100] {1 - exp(-x)};
						\addlegendentry{$u_C(t)/U_0$}
						\addplot[red, thick, domain=0:5, samples=100] {exp(-x)};
						\addlegendentry{$i_C(t)\cdot R/U_0$}
					\end{axis}
				\end{tikzpicture}
			\end{minipage}
			\caption{$u$-$i$-Diagramm: Laden des Kondensators}
		\end{subfigure}
		\caption{RC-Ladeschaltung: $u_C$ steigt asymptotisch auf $U_0$, während $i_C$ exponentiell abfällt. Beide Kurven bleiben stets $\geq 0$.}
		\label{fig:rc_laden}
	\end{figure}
	
	\subsection{Physikalisch-mathematische Konsistenz}
	
	Da die Ladung $Q = C \cdot U$ und Spannung $U$ beim Kondensator proportional verknüpft sind, folgt:
	\begin{equation}
		W_C = Q \cdot U / 2 = (C \cdot U) \cdot U / 2 = \frac{1}{2} C \cdot U^2 \geq 0.
	\end{equation}
	
	Das $u$-$i$-Diagramm verdeutlicht: Strom und Spannung sind stets nicht-negativ beim Ladevorgang. Eine Umpolung der Spannung würde dasselbe Ergebnis liefern: $(-U)^2 = U^2$. Der Kondensator speichert stets Energie, unabhängig von der Polarität.
	
	% ============================================================
	\section{Die elektrische Spule (Induktivität)}
	% ============================================================
	
	\subsection{Physikalische Beschreibung und Schaltung}
	
	Die in einer Spule gespeicherte magnetische Energie beträgt:
	\begin{equation}
		W_L = \frac{1}{2} L \cdot I^2 \geq 0
		\label{eq:spule_energie}
	\end{equation}
	
	Die induzierte Spannung und der Stromaufbau beim Einschalten lauten:
	\begin{equation}
		u_L(t) = U_0 \cdot e^{-Rt/L}, \qquad
		i_L(t) = \frac{U_0}{R}\left(1 - e^{-Rt/L}\right)
	\end{equation}
	
	\begin{figure}[h]
		\centering
		\begin{subfigure}[b]{0.48\textwidth}
			\centering
			\begin{minipage}[c][6.5cm][c]{\linewidth}
				\centering
				\begin{circuitikz}[scale=1.0, transform shape]
					\draw (0,0) to[battery1, l=$U_0$, invert] (0,3)
					to[R, l=$R$] (3,3)
					to[L, l=$L$, v=$u_L$] (3,0)
					to[short] (0,0);
					% \draw[->, thick] (1.5,3.2) -- (2.5,3.2) node[midway, above]{$i_L$};
				\end{circuitikz}
			\end{minipage}
			\caption{RL-Einschaltschaltung}
		\end{subfigure}
		\hfill
		\begin{subfigure}[b]{0.48\textwidth}
			\centering
			\begin{minipage}[c][6.5cm][c]{\linewidth}
				\centering
				\begin{tikzpicture}
					\begin{axis}[
						% xlabel={Zeit $t / \tau_L$},
						ylabel={Normierte Größe},
						xmin=0, xmax=5,
						ymin=0, ymax=1.1,
						width=7cm, height=5cm,
						grid=major,
						legend pos=north east,
						]
						\addplot[red, thick, domain=0:5, samples=100] {1 - exp(-x)};
						\addlegendentry{$i_L(t) \cdot R/U_0$}
						\addplot[blue, thick, domain=0:5, samples=100] {exp(-x)};
						\addlegendentry{$u_L(t)/U_0$}
					\end{axis}
				\end{tikzpicture}
			\end{minipage}
			\caption{$u$-$i$-Diagramm: Einschalten der Spule}
		\end{subfigure}
		\caption{RL-Einschaltschaltung: $i_L$ steigt asymptotisch an, $u_L$ fällt exponentiell. Beide Größen bleiben $\geq 0$.}
		\label{fig:rl_einschalten}
	\end{figure}
	
	\subsection{Physikalisch-mathematische Konsistenz}
	
	Der magnetische Fluss $\Phi = L \cdot I$ und der Strom $I$ sind proportional verknüpft. Damit:
	\begin{equation}
		W_L = \int_0^I L \cdot I' \, dI' = \frac{1}{2} L \cdot I^2 \geq 0.
	\end{equation}
	
	Die Spule ist das duale Element zum Kondensator: Während der Kondensator Energie im elektrischen Feld speichert ($W_C \sim U^2$), speichert die Spule Energie im magnetischen Feld ($W_L \sim I^2$). In beiden Fällen erzwingt die Proportionalität der gespeicherten Größen die quadratische Form -- und damit die Nicht-Negativität der Energie.
	
	% ============================================================
	\section{Das RC-Entladeglied}
	% ============================================================
	
	\subsection{Physikalische Beschreibung und Schaltung}
	
	Beim Entladen eines Kondensators über einen Widerstand ergibt sich:
	\begin{equation}
		u_C(t) = U_0 \cdot e^{-t/\tau}, \qquad \tau = RC
	\end{equation}
	\begin{equation}
		i(t) = -\frac{U_0}{R} \cdot e^{-t/\tau}
	\end{equation}
	
	Der Strom ist während der Entladung negativ (er fließt entgegengesetzt zur Laderichtung), die Spannung bleibt positiv und fällt gegen null.
	
	\begin{figure}[h]
		\centering
		\begin{subfigure}[b]{0.48\textwidth}
			\centering
			\begin{minipage}[c][6.5cm][c]{\linewidth}
				\centering
				\begin{circuitikz}[scale=1.0, transform shape]
					% Charged capacitor discharging through R
					\draw (0,0) to[C, l=$C$, v=$u_C$] (0,3)
					to[short] (3,3)
					to[R, l=$R$] (3,0)
					to[short] (0,0);
					% \draw[->, thick, red] (1.5,3.2) -- (2.5,3.2) node[midway, above]{$i(t)$};
					\node at (1.5,-0.5) {Entladeschaltung};
				\end{circuitikz}
			\end{minipage}
			\caption{RC-Entladeschaltung}
		\end{subfigure}
		\hfill
		\begin{subfigure}[b]{0.48\textwidth}
			\centering
			\begin{minipage}[c][6.5cm][c]{\linewidth}
				\centering
				\begin{tikzpicture}
					\begin{axis}[
						% xlabel={Zeit $t / \tau$},
						ylabel={Normierte Größe},
						xmin=0, xmax=5,
						ymin=-1.1, ymax=1.1,
						width=7cm, height=5cm,
						grid=major,
						legend pos=north east,
						]
						\addplot[blue, thick, domain=0:5, samples=100] {exp(-x)};
						\addlegendentry{$u_C(t)/U_0 \geq 0$}
						\addplot[red, thick, domain=0:5, samples=100] {-exp(-x)};
						\addlegendentry{$i(t) \cdot R/U_0 \leq 0$}
						\addplot[dashed, gray, domain=0:5] {0};
					\end{axis}
				\end{tikzpicture}
			\end{minipage}
			\caption{$u$-$i$-Diagramm: Entladevorgang}
		\end{subfigure}
		\caption{RC-Entladeschaltung: Spannung fällt positiv ab, Strom ist negativ (Entladung). Die Energie $W = \frac{1}{2}Cu_C^2$ ist stets $\geq 0$.}
		\label{fig:rc_entladen}
	\end{figure}
	
	\subsection{Physikalisch-mathematische Konsistenz}
	
	Obwohl der Entladestrom negativ ist, ist die gespeicherte Energie im Kondensator stets positiv:
	\begin{equation}
		W_C(t) = \frac{1}{2} C \cdot u_C^2(t) = \frac{1}{2} C \cdot U_0^2 \cdot e^{-2t/\tau} \geq 0.
	\end{equation}
	
	Das $u$-$i$-Diagramm illustriert die asymmetrische Situation: Spannung und Strom haben unterschiedliche Vorzeichen während der Entladung -- was physikalisch korrekt ist, da die Energie aus dem Kondensator in den Widerstand fließt. Die dissipierte Leistung im Widerstand:
	\begin{equation}
		P_R(t) = R \cdot i^2(t) = \frac{U_0^2}{R} e^{-2t/\tau} \geq 0
	\end{equation}
	ist erneut nicht-negativ, da das Quadrat des Stromes stets positiv ist.
	
	% ============================================================
	\section{Der Halbwellengleichrichter (Einweggleichrichter)}
	% ============================================================
	
	\subsection{Physikalische Beschreibung und Schaltung}
	
	Ein einzelne Diode in Reihe mit einem Lastwiderstand bildet den einfachsten Gleichrichter. Nur die positive Halbwelle der Wechselspannung wird durchgelassen.
	
	\begin{figure}[h]
		\centering
		\begin{subfigure}[b]{0.48\textwidth}
			\centering
			\begin{minipage}[c][6.5cm][c]{\linewidth}
				\centering
				\begin{circuitikz}[scale=1.0, transform shape]
					\draw (0,0) to[sV, l=$u_{in}$] (0,3)
					to[D, l=D] (3,3)
					to[R, l=$R_L$, v=$u_{out}$] (3,0)
					to[short] (0,0);
					% \draw[->, thick] (1.5,3.2) -- (2.5,3.2) node[midway, above]{$i$};
				\end{circuitikz}
			\end{minipage}
			\caption{Halbwellengleichrichter}
		\end{subfigure}
		\hfill
		\begin{subfigure}[b]{0.48\textwidth}
			\centering
			\begin{minipage}[c][6.5cm][c]{\linewidth}
				\centering
				\begin{tikzpicture}
					\begin{axis}[
						% xlabel={Zeit $\omega t$},
						ylabel={Normierte Spannung},
						xmin=0, xmax=6.28,
						ymin=-1.1, ymax=1.1,
						width=7cm, height=5cm,
						grid=major,
						legend pos=north east,
						xtick={0,1.57,3.14,4.71,6.28},
						xticklabels={$0$,$\frac{\pi}{2}$,$\pi$,$\frac{3\pi}{2}$,$2\pi$}
						]
						\addplot[blue, thin, dashed, domain=0:6.28, samples=100] {sin(deg(x))};
						\addlegendentry{$u_{in}$}
						\addplot[red, thick, domain=0:6.28, samples=200]
						{(sin(deg(x)) > 0) ? sin(deg(x)) : 0};
						\addlegendentry{$u_{out} \geq 0$}
					\end{axis}
				\end{tikzpicture}
			\end{minipage}
			\caption{$u$-$i$-Diagramm: Halbwellengleichrichtung}
		\end{subfigure}
		\caption{Halbwellengleichrichter: Die negative Halbwelle wird durch die Diode gesperrt. Die Ausgangsspannung entspricht mathematisch der Funktion $u_{out} = \max(u_{in},\,0)$.}
		\label{fig:halbwelle}
	\end{figure}
	
	\subsection{Physikalisch-mathematische Konsistenz}
	
	Die mathematische Entsprechung des Halbwellengleichrichters ist die Klammerfunktion:
	\begin{equation}
		u_{out}(t) = \max\bigl(u_{in}(t),\; 0\bigr) = \frac{u_{in}(t) + |u_{in}(t)|}{2}.
	\end{equation}
	
	Die Diode implementiert diese Funktion physikalisch. Auch hier gilt: Eine rein mathematische Betrachtung des Sinussignals ohne Berücksichtigung der Diode würde negative Ausgangsspannungen vorhersagen -- die physikalische Struktur schließt dies aus.
	
	% ============================================================
	\section{Der Brückengleichrichter (Vollwellengleichrichter)}
	% ============================================================
	
	\subsection{Physikalische Beschreibung und Schaltung}
	
	Vier Dioden in Brückenschaltung (Graetz-Brücke) ermöglichen die Gleichrichtung beider Halbwellen.
	
	\begin{figure}[h]
		\centering
		\begin{subfigure}[b]{0.48\textwidth}
			\centering
			\begin{minipage}[c][6.5cm][c]{\linewidth}
				\centering
				\begin{circuitikz}[scale=0.9, transform shape]
					% Graetz bridge
					\draw (2,0) to[D, l=$D_3$] (0,2);
					\draw (0,2) to[D, l=$D_1$] (2,4);
					\draw (2,0) to[D, l=$D_4$, invert] (4,2);
					\draw (4,2) to[D, l=$D_2$, invert] (2,4);
					% AC source
					\draw (0,2) to[short] (-1,2) to[sV, l=$u_{in}$] (-1,2) to[short] (-1,2);
					\draw (4,2) to[short] (5,2);
					\draw (-1,0.5) to[sV, l=$u_{in}$] (-1,3.5);
					\draw (-1,0.5) to[short] (2,0);
					\draw (-1,3.5) to[short] (2,4);
					% Load
					\draw (5,2) to[short] (5,3.5) to[R, l=$R_L$, v=$u_{out}$] (5,0.5) to[short] (2,0);
					\draw (5,2) to[short] (2,4);
				\end{circuitikz}
			\end{minipage}
			\caption{Graetz-Brückenschaltung}
		\end{subfigure}
		\hfill
		\begin{subfigure}[b]{0.48\textwidth}
			\centering
			\begin{minipage}[c][6.5cm][c]{\linewidth}
				\centering
				\begin{tikzpicture}
					\begin{axis}[
						% xlabel={Zeit $\omega t$},
						ylabel={Normierte Spannung},
						xmin=0, xmax=6.28,
						ymin=-1.1, ymax=1.1,
						width=7cm, height=5cm,
						grid=major,
						legend pos=north east,
						xtick={0,1.57,3.14,4.71,6.28},
						xticklabels={$0$,$\frac{\pi}{2}$,$\pi$,$\frac{3\pi}{2}$,$2\pi$}
						]
						\addplot[blue, thin, dashed, domain=0:6.28, samples=100] {sin(deg(x))};
						\addlegendentry{$u_{in}$}
						\addplot[red, thick, domain=0:6.28, samples=200] {abs(sin(deg(x)))};
						\addlegendentry{$u_{out} = |u_{in}|$}
					\end{axis}
				\end{tikzpicture}
			\end{minipage}
			\caption{$u$-$i$-Diagramm: Vollwellengleichrichtung}
		\end{subfigure}
		\caption{Graetz-Brückengleichrichter: Beide Halbwellen werden gleichgerichtet. Die Ausgangsspannung entspricht mathematisch der Betragsfunktion $|u_{in}|$.}
		\label{fig:bruecke}
	\end{figure}
	
	\subsection{Physikalisch-mathematische Konsistenz}
	
	Der Brückengleichrichter implementiert die mathematische Betragsfunktion:
	\begin{equation}
		u_{out}(t) = \bigl|u_{in}(t)\bigr| = \hat{U} \cdot |\sin(\omega t)| \geq 0.
	\end{equation}
	
	Dies ist ein besonders elegantes Beispiel der Einheit von Physik und Mathematik: Die Betragsfunktion, die in der Mathematik als einfache Vorzeichenentfernung definiert ist, wird durch eine spezifische Anordnung von vier Halbleitern physikalisch realisiert. Der mittlere Gleichspannungswert beträgt:
	\begin{equation}
		\bar{U} = \frac{2\hat{U}}{\pi} \approx 0{,}637 \cdot \hat{U}.
	\end{equation}
	
	% ============================================================
	\section{Der Spannungsteiler}
	% ============================================================
	
	\subsection{Physikalische Beschreibung und Schaltung}
	
	Zwei Widerstände in Reihe bilden einen Spannungsteiler. Die Ausgangsspannung verhält sich zur Eingangsspannung wie das Widerstandsverhältnis:
	\begin{equation}
		U_{out} = U_{in} \cdot \frac{R_2}{R_1 + R_2}
	\end{equation}
	
	\begin{figure}[h]
		\centering
		\begin{subfigure}[b]{0.48\textwidth}
			\centering
			\begin{minipage}[c][6.5cm][c]{\linewidth}
				\centering
				\begin{circuitikz}[scale=1.0, transform shape]
					\draw (0,4) to[short, -o] (-0.5,4) node[left]{$+$}
					(0,4) to[R, l=$R_1$] (0,2)
					to[R, l=$R_2$, v=$U_{out}$] (0,0)
					to[short, -o] (-0.5,0) node[left]{$-$};
					\draw (-0.5,4) to[open, v=$U_{in}$] (-0.5,0);
					\draw (0,2) to[short, -o] (1,2) node[right]{$U_{out}$};
					\draw (0,0) to[short, -o] (1,0) node[right]{GND};
				\end{circuitikz}
			\end{minipage}
			\caption{Spannungsteiler}
		\end{subfigure}
		\hfill
		\begin{subfigure}[b]{0.48\textwidth}
			\centering
			\begin{minipage}[c][6.5cm][c]{\linewidth}
				\centering
				\begin{tikzpicture}
					\begin{axis}[
						% xlabel={$R_2 / (R_1 + R_2)$},
						ylabel={$U_{out} / U_{in}$},
						xmin=0, xmax=1,
						ymin=0, ymax=1.05,
						width=7cm, height=5cm,
						grid=major,
						legend pos=north west,
						]
						\addplot[blue, thick, domain=0:1, samples=50] {x};
						\addlegendentry{$U_{out}/U_{in} = k \in [0,1]$}
						\addplot[red, dashed, domain=0:1] {0.5};
						\node at (axis cs:0.5,0.53) {$k = 0{,}5$};
					\end{axis}
				\end{tikzpicture}
			\end{minipage}
			\caption{Teilungsverhältnis: stets $0 \leq U_{out} \leq U_{in}$}
		\end{subfigure}
		\caption{Spannungsteiler: Das Teilungsverhältnis liegt stets im Bereich $[0,1]$, sofern $R_1, R_2 > 0$. Die Mathematik garantiert $0 \leq U_{out} \leq U_{in}$.}
		\label{fig:spannungsteiler}
	\end{figure}
	
	\subsection{Physikalisch-mathematische Konsistenz}
	
	Das Teilungsverhältnis $k = R_2/(R_1+R_2)$ liegt für alle positiven Widerstände $R_1, R_2 > 0$ stets im Intervall $(0,1)$. Die Mathematik garantiert: $0 < U_{out} < U_{in}$. Ein Spannungsteiler kann nie eine höhere Spannung ausgeben als er eingangsseitig erhält -- physikalisch selbstverständlich, mathematisch durch die Ungleichung $R_2 < R_1 + R_2$ präzise ausgedrückt.
	
	% ============================================================
	\section{Der ideale Transformator}
	% ============================================================
	
	\subsection{Physikalische Beschreibung und Schaltung}
	
	Ein idealer Transformator mit Übersetzungsverhältnis $\ddot{u} = N_1 / N_2$ transformiert:
	\begin{equation}
		U_2 = \frac{U_1}{\ddot{u}}, \qquad I_2 = \ddot{u} \cdot I_1, \qquad P_1 = P_2
		\label{eq:trafo}
	\end{equation}
	
	\begin{figure}[h]
		\centering
		\begin{subfigure}[b]{0.48\textwidth}
			\centering
			\begin{minipage}[c][6.5cm][c]{\linewidth}
				\centering
				\begin{circuitikz}[scale=1.0, transform shape]
					\draw (0,0) to[sV, l=$U_1$] (0,3)
					to[short] (1.5,3);
					\draw (1.5,0) to[short] (0,0);
					\draw (1.5,0) to[transformer core] (3.5,0);
					\node at (2.5,1.5) {$\ddot{u}{=}N_1/N_2$};
					\draw (3.5,3) to[short] (5,3)
					to[R, l=$R_L$, v=$U_2$] (5,0)
					to[short] (3.5,0);
					\node[above] at (0.75,3) {$N_1$};
					\node[above] at (4.25,3) {$N_2$};
				\end{circuitikz}
			\end{minipage}
			\caption{Idealer Transformator mit Last}
		\end{subfigure}
		\hfill
		\begin{subfigure}[b]{0.48\textwidth}
			\centering
			\begin{minipage}[c][6.5cm][c]{\linewidth}
				\centering
				\begin{tikzpicture}
					\begin{axis}[
						% xlabel={Zeit $\omega t$},
						ylabel={Normierte Größe},
						xmin=0, xmax=6.28,
						ymin=-1.6, ymax=1.6,
						width=7cm, height=5cm,
						grid=major,
						legend pos=north east,
						xtick={0,3.14,6.28},
						xticklabels={$0$,$\pi$,$2\pi$}
						]
						\addplot[blue, thick, domain=0:6.28, samples=100] {sin(deg(x))};
						\addlegendentry{$u_1(t)/\hat{U}_1$}
						\addplot[red, thick, dashed, domain=0:6.28, samples=100] {1.5*sin(deg(x))};
						\addlegendentry{$u_2(t)/\hat{U}_1$ ($\ddot{u}{<}1$)}
						\addplot[green!60!black, thick, dotted, domain=0:6.28, samples=100] {(1/1.5)*sin(deg(x))};
						\addlegendentry{$i_2 \cdot R / \hat{U}_1$}
					\end{axis}
				\end{tikzpicture}
			\end{minipage}
			\caption{$u$-$i$-Diagramm: Leistungserhaltung}
		\end{subfigure}
		\caption{Idealer Transformator: Obwohl $U_2 \neq U_1$, gilt stets $P_1 = P_2 = U_1 I_1 = U_2 I_2$. Das Produkt aus Spannung und Strom ist invariant.}
		\label{fig:transformator}
	\end{figure}
	
	\subsection{Physikalisch-mathematische Konsistenz}
	
	\begin{equation}
		P_2 = U_2 \cdot I_2 = \frac{U_1}{\ddot{u}} \cdot (\ddot{u} \cdot I_1) = U_1 \cdot I_1 = P_1.
	\end{equation}
	
	Das $u$-$i$-Diagramm zeigt: Während die Amplituden von Spannung und Strom auf der Sekundärseite verändert werden, bleibt das Produkt -- die Leistung -- konstant. Die Mathematik des idealen Transformators macht das Prinzip der Energieerhaltung unmittelbar sichtbar.
	
	% ============================================================
	\section{Der Schwingkreis (LC-Kreis)}
	% ============================================================
	
	\subsection{Physikalische Beschreibung und Schaltung}
	
	Ein Parallelschwingkreis aus Spule $L$ und Kondensator $C$ zeigt ungedämpfte harmonische Schwingungen. Die Gesamtenergie bleibt erhalten:
	\begin{equation}
		W_{ges} = \frac{1}{2}L i_L^2 + \frac{1}{2}C u_C^2 = \text{const}
	\end{equation}
	
	\begin{figure}[h]
		\centering
		\begin{subfigure}[b]{0.48\textwidth}
			\centering
			\begin{minipage}[c][6.5cm][c]{\linewidth}
				\centering
				\begin{circuitikz}[scale=1.0, transform shape]
					% Parallel LC circuit
					\draw (0,3) to[short] (4,3);
					\draw (0,0) to[short] (4,0);
					\draw (0,0) to[C, l=$C$, v=$u_C$] (0,3);
					\draw (2,0) to[L, l=$L$] (2,3);
					% \draw[->, thick] (3.5,2.5) -- (3.5,0.5) node[midway, right]{$i_L$};
					% \node at (2,-0.6) {Parallelschwingkreis};
				\end{circuitikz}
			\end{minipage}
			\caption{LC-Parallelschwingkreis}
		\end{subfigure}
		\hfill
		\begin{subfigure}[b]{0.48\textwidth}
			\centering
			\begin{minipage}[c][6.5cm][c]{\linewidth}
				\centering
				\begin{tikzpicture}
					\begin{axis}[
						% xlabel={Zeit $\omega_0 t$},
						ylabel={Normierte Größe},
						xmin=0, xmax=6.28,
						ymin=-1.1, ymax=1.1,
						width=7cm, height=5cm,
						grid=major,
						legend pos=north east,
						xtick={0,1.57,3.14,4.71,6.28},
						xticklabels={$0$,$\frac{\pi}{2}$,$\pi$,$\frac{3\pi}{2}$,$2\pi$}
						]
						\addplot[blue, thick, domain=0:6.28, samples=100] {cos(deg(x))};
						\addlegendentry{$u_C(t)/\hat{U}$}
						\addplot[red, thick, domain=0:6.28, samples=100] {sin(deg(x))};
						\addlegendentry{$i_L(t)/\hat{I}$}
						\addplot[green!60!black, thick, domain=0:6.28, samples=100]
						{cos(deg(x))^2 + sin(deg(x))^2};
						\addlegendentry{$W_{ges} = 1$ (const.)}
					\end{axis}
				\end{tikzpicture}
			\end{minipage}
			\caption{$u$-$i$-Diagramm: Energieerhaltung}
		\end{subfigure}
		\caption{LC-Schwingkreis: Spannung und Strom schwingen um $90^\circ$ phasenverschoben. Die Gesamtenergie bleibt konstant: $\cos^2(\omega_0 t) + \sin^2(\omega_0 t) = 1$.}
		\label{fig:schwingkreis}
	\end{figure}
	
	\subsection{Physikalisch-mathematische Konsistenz}
	
	Die Differentialgleichung des LC-Kreises lautet:
	\begin{equation}
		L \ddot{i}_L + \frac{i_L}{C} = 0 \quad \Rightarrow \quad \ddot{i}_L + \omega_0^2 i_L = 0, \quad \omega_0 = \frac{1}{\sqrt{LC}}.
	\end{equation}
	
	Die Lösung ist $i_L(t) = \hat{I}\sin(\omega_0 t)$, $u_C(t) = \hat{U}\cos(\omega_0 t)$. Die Gesamtenergie:
	\begin{equation}
		W_{ges} = \frac{1}{2}L\hat{I}^2\sin^2(\omega_0 t) + \frac{1}{2}C\hat{U}^2\cos^2(\omega_0 t) = \frac{1}{2}L\hat{I}^2,
	\end{equation}
	da $\hat{I} = \hat{U}/(\omega_0 L)$. Die Summe $\sin^2 + \cos^2 = 1$ -- der trigonometrische Pythagoras -- ist exakt das mathematische Abbild der physikalischen Energieerhaltung im verlustlosen Schwingkreis.
	
	% ============================================================
	\section{Das RLC-Serienglied: Gedämpfte Schwingung}
	% ============================================================
	
	\subsection{Physikalische Beschreibung und Schaltung}
	
	Ein Reihenschwingkreis aus $R$, $L$, $C$ zeigt je nach Dämpfung unterschiedliches Verhalten:
	\begin{equation}
		L \ddot{u}_C + R \dot{u}_C + \frac{u_C}{C} = 0
	\end{equation}
	
	\begin{figure}[h]
		\centering
		\begin{subfigure}[b]{0.48\textwidth}
			\centering
			\begin{minipage}[c][6.5cm][c]{\linewidth}
				\centering
				\begin{circuitikz}[scale=1.0, transform shape]
					\draw (0,0) to[battery1, l=$U_0$, invert] (0,3)
					to[R, l=$R$] (2,3)
					to[L, l=$L$] (4,3)
					to[C, l=$C$, v=$u_C$] (4,0)
					to[short] (0,0);
					% \draw[->, thick] (1,3.2) -- (3,3.2) node[midway, above]{$i(t)$};
				\end{circuitikz}
			\end{minipage}
			\caption{RLC-Serienschwingkreis}
		\end{subfigure}
		\hfill
		\begin{subfigure}[b]{0.48\textwidth}
			\centering
			\begin{minipage}[c][6.5cm][c]{\linewidth}
				\centering
				\begin{tikzpicture}
					\begin{axis}[
						% xlabel={Zeit $\omega_0 t$},
						ylabel={Normierte Spannung $u_C/U_0$},
						xmin=0, xmax=20,
						ymin=-0.5, ymax=1.5,
						width=7cm, height=5cm,
						grid=major,
						legend pos=north east,
						]
						% Underdamped: u_C = 1 - e^(-delta*t)*cos(omega_d*t) + ...
						\addplot[blue, thick, domain=0:20, samples=200]
						{1 - exp(-0.2*x)*(cos(deg(x)) + 0.2*sin(deg(x)))};
						\addlegendentry{$u_C$ (schwach ged.)}
						\addplot[red, thick, domain=0:20, samples=100]
						{1 - exp(-x)};
						\addlegendentry{$u_C$ (krit. ged.)}
					\end{axis}
				\end{tikzpicture}
			\end{minipage}
			\caption{$u_C$-Zeitverlauf: Unter- und Kriechfall}
		\end{subfigure}
		\caption{RLC-Serienschwingkreis: Schwache Dämpfung ($R$ klein) erzeugt gedämpfte Schwingung (blau). Kritische Dämpfung (rot) führt zum schnellsten aperiodischen Einlauf.}
		\label{fig:rlc}
	\end{figure}
	
	\subsection{Physikalisch-mathematische Konsistenz}
	
	Der Dämpfungsgrad $\zeta = R/(2\sqrt{L/C})$ bestimmt das Verhalten vollständig:
	\begin{itemize}
		\item $\zeta < 1$: Schwingfall -- die Wurzeln der charakteristischen Gleichung sind komplex
		\item $\zeta = 1$: Kriechfall -- doppelte reelle Wurzel
		\item $\zeta > 1$: Überdämpfung -- zwei verschiedene reelle Wurzeln
	\end{itemize}
	
	In allen Fällen gilt: Die Gesamtenergie des Systems nimmt monoton ab (Dissipation im Widerstand), und die Kondensatorspannung nähert sich asymptotisch dem Endwert $U_0$. Die Mathematik -- die Theorie der linearen Differentialgleichungen 2. Ordnung -- beschreibt das physikalische Verhalten vollständig und lückenlos.

	\subsection{Der Schwerpunkt des Schwingkreises: Die Spule}

	\textbf{Grundsatz.} \textit{Physikalische Strukturen sind in ihrem Schwerpunkt am leichtesten zu lösen.} Der Schwerpunkt einer physikalischen Struktur ist dasjenige Bauelement, von dem aus die Differentialgleichung ihre natürlichste, klarste Form annimmt -- jene Form, in der jeder Koeffizient eine unmittelbar physikalisch interpretierbare Bedeutung trägt.

	Für den RCL-Schwingkreis lässt sich die Differentialgleichung prinzipiell auf zwei Wegen aufstellen: über den Strom durch die Spule oder über die Spannung am Kondensator. Beide Wege führen formal zur gleichen charakteristischen Gleichung -- aber nur einer von beiden offenbart die Struktur unmittelbar, ohne zusätzliche Normierungsschritte. Dieser Weg verläuft über die Spule.

	\textbf{Formulierung über die Spule (Strom als Variable).}
	Für den RCL-Serienkreis liefert die Kirchhoffsche Maschenregel:
	\begin{equation}
		u_L + u_R + u_C = 0
		\quad \Longrightarrow \quad
		L\,\frac{di}{dt} + R\,i + \frac{1}{C}\int i\,dt = 0.
	\end{equation}
	Einmalige Differentiation nach der Zeit ergibt direkt:
	\begin{equation}
		L\,\ddot{i} + R\,\dot{i} + \frac{1}{C}\,i = 0.
		\label{eq:rlc_spule}
	\end{equation}
	Division durch $L$ bringt die Gleichung sofort in die Standardform des gedämpften harmonischen Oszillators:
	\begin{equation}
		\ddot{i} + \underbrace{\frac{R}{L}}_{2\delta}\,\dot{i} + \underbrace{\frac{1}{LC}}_{\omega_0^2}\,i = 0,
		\qquad \delta = \frac{R}{2L},\quad \omega_0 = \frac{1}{\sqrt{LC}}.
		\label{eq:rlc_standardform}
	\end{equation}
	Der Dämpfungskoeffizient $\delta$ und die Eigenfrequenz $\omega_0$ lassen sich unmittelbar aus den drei Bauelementwerten ablesen. Kein weiterer Umformungsschritt ist erforderlich.

	\textbf{Formulierung über den Kondensator (Spannung als Variable).}
	Wählt man stattdessen $u_C$ als Zustandsgröße und setzt $i = C\,\dot{u}_C$, so folgt:
	\begin{equation}
		LC\,\ddot{u}_C + RC\,\dot{u}_C + u_C = 0.
		\label{eq:rlc_kondensator}
	\end{equation}
	Hier steht vor dem höchsten Differentialterm das \textit{Produkt} $LC$ zweier verschiedener Bauelemente -- kein einzelner, physikalisch eigenständiger Parameter. Erst nach Division durch $LC$ entsteht dieselbe Standardform wie in \eqref{eq:rlc_standardform}. Die Struktur ist algebraisch vorhanden, aber nicht unmittelbar sichtbar.

	\textbf{Die Spule als Trägheitselement.}
	Der Grund für diese Asymmetrie liegt in der physikalischen Rolle der Spule. Genau wie die Masse $m$ in der Newtonschen Gleichung
	\begin{equation}
		m\,\ddot{x} + d\,\dot{x} + k\,x = 0
	\end{equation}
	dem mechanischen Schwinger seine Trägheit verleiht, so verleiht die Induktivität $L$ dem Schwingkreis seine \textit{elektrische Trägheit}: Sie widersetzt sich jeder Änderung des Stromes. Die Analogie ist vollständig:
	\begin{equation}
		\underbrace{L}_{\displaystyle\leftrightarrow\, m}\,\ddot{i}
		+ \underbrace{R}_{\displaystyle\leftrightarrow\, d}\,\dot{i}
		+ \underbrace{\frac{1}{C}}_{\displaystyle\leftrightarrow\, k}\,i = 0.
	\end{equation}
	Das Produkt $L\,i$ ist der magnetische Fluss $\Phi = L\,i$ -- das elektromagnetische Analogon des mechanischen Impulses $p = m\,v$. Die Spule \textit{speichert den Impuls} des Schwingkreises; der Kondensator speichert die potenzielle Energie ($W_C = \frac{1}{2}Cu_C^2$, analog zur Federenergie $\frac{1}{2}kx^2$).

	Die Differentialgleichung \eqref{eq:rlc_spule}, formuliert über den Stromfluss durch die Spule, ist daher diejenige, von der aus die Dynamik des Schwingkreises am direktesten und klarsten beschrieben wird. \textbf{Die Spule ist} -- im Sinne des eingangs formulierten Grundsatzes -- \textbf{der Schwerpunkt des RCL-Schwingkreises}.
	
	
	\section{Der Schwerpunkt-Grundsatz: Anwendung auf alle Spulenschaltungen}
	\label{sec:schwerpunkt_spule_allgemein}
	
	\subsection{Formale Präzisierung des Grundsatzes}
	
	Der in Abschnitt~11.3 formulierte Grundsatz lautet:
	
	\begin{quote}
		\textit{Physikalische Strukturen sind in ihrem Schwerpunkt am leichtesten zu lösen.
			Der Schwerpunkt ist dasjenige Bauelement, von dem aus die Differentialgleichung
			ihre natürlichste, klarste Form annimmt -- jene Form, in der jeder Koeffizient
			eine unmittelbar physikalisch interpretierbare Bedeutung trägt.}
	\end{quote}
	
	Bevor dieser Grundsatz auf alle Spulenschaltungen der Arbeit angewendet wird,
	soll er zunächst mathematisch präzise formuliert werden. Gegeben sei eine
	lineare Schaltung mit den Zustandsgrößen $x_1, x_2, \ldots, x_n$ (Ströme
	durch Spulen, Spannungen über Kondensatoren). Die allgemeine Zustandsraum\-form
	lautet:
	\begin{equation}
		\dot{\mathbf{x}} = A\,\mathbf{x} + B\,\mathbf{u},
		\label{eq:zustandsraum}
	\end{equation}
	wobei $A$ die Systemmatrix, $B$ die Eingangsmatrix und $\mathbf{u}$ die
	Erregung ist.
	
	\textbf{Definition (Schwerpunkt-Variable).}
	Eine Zustandsgröße $x_k$ heißt \textit{Schwerpunkt-Variable} der Schaltung,
	wenn die aus ihr abgeleitete skalare Differentialgleichung (nach algebraischer
	Elimination aller übrigen Zustände) die Eigenschaft besitzt, dass der
	Leitkoeffizient (Koeffizient der höchsten Ableitung) eine einzelne,
	physikalisch eigenständige Bauelementegröße ist -- und nicht ein Produkt oder
	Quotient mehrerer Bauelementegrößen.
	
	Formal: Sei $p(\mathfrak{D})$ das charakteristische Polynom in der
	Differentialoperator-Darstellung. Die Differentialgleichung für $x_k$ lautet:
	\begin{equation}
		p(\mathfrak{D})\,x_k = f(t).
	\end{equation}
	$x_k$ ist Schwerpunkt-Variable, wenn $p(\mathfrak{D})$ die Form
	\begin{equation}
		p(\mathfrak{D}) = \alpha_k\,\mathfrak{D}^n + a_{n-1}\,\mathfrak{D}^{n-1}
		+ \cdots + a_0
	\end{equation}
	besitzt, wobei $\alpha_k$ die physikalische Größe \emph{genau eines} Bauelements ist.
	
	Die Spule ist in allen nachfolgend betrachteten Schaltungen die
	Schwerpunkt-Variable, weil ihre Induktivität $L$ stets und direkt als
	Leitkoeffizient $\alpha_k = L$ erscheint. Dies ist keine Zufälligkeit,
	sondern Konsequenz ihrer physikalischen Rolle als Trägheitselement: $L$
	ist das elektrische Analogon der Masse $m$, und ebenso wie die Masse in
	der Newtonschen Mechanik als Koeffizient der höchsten Zeitableitung
	($m\ddot{x}$) erscheint, erscheint $L$ als Koeffizient von $\ddot{i}$.
	
	\subsection{Die RL-Einschaltschaltung: Die Spule als einziger Energieträger}
	\label{sec:schwerpunkt_rl}
	
	\subsubsection*{Die Schaltung}
	
	Die einfachste Schaltung mit einer Spule ist das RL-Serienglied an einer
	Gleichspannungsquelle $U_0$ (vgl.\ Abschnitt~9 dieser Arbeit). Die
	Kirchhoffsche Maschengleichung lautet:
	\begin{equation}
		U_0 = u_R + u_L = R\,i + L\,\frac{di}{dt}.
		\label{eq:rl_masche}
	\end{equation}
	
	\subsubsection*{Warum die Spule der Schwerpunkt ist}
	
	Gleichung~\eqref{eq:rl_masche} ist \emph{bereits} in der Schwerpunkt-Form
	geschrieben: Der Leitkoeffizient des höchsten Differentialterms ist allein
	$L$. Umstellen nach $\frac{di}{dt}$ ergibt sofort:
	\begin{equation}
		\frac{di}{dt} = \frac{U_0 - R\,i}{L}.
		\label{eq:rl_dgl}
	\end{equation}
	Division durch $L$ liefert die normierte Standardform:
	\begin{equation}
		\dot{i} + \underbrace{\frac{R}{L}}_{\displaystyle 1/\tau_L}\,i
		= \frac{U_0}{L},
		\qquad \tau_L = \frac{L}{R}.
		\label{eq:rl_standardform}
	\end{equation}
	Die Zeitkonstante $\tau_L = L/R$ ist das einzige dynamische
	Strukturparameter der Schaltung -- und sie lässt sich unmittelbar aus der
	Spule ablesen: $L$ im Zähler (Trägheit), $R$ im Nenner (Verlust).
	
	\textbf{Vergleich: Formulierung über den Widerstand.}
	Alternativ könnte man die Spannung $u_R = R\,i$ als primäre Variable wählen.
	Da $u_R = U_0 - u_L$ und $u_L = L\,\dot{i} = (L/R)\,\dot{u}_R$, folgt:
	\begin{equation}
		\frac{L}{R}\,\dot{u}_R + u_R = U_0.
		\label{eq:rl_widerstand}
	\end{equation}
	Diese Gleichung ist formal identisch zu \eqref{eq:rl_standardform}, aber
	der Leitkoeffizient ist nun $L/R$ -- das \emph{Verhältnis} zweier
	Bauelementegrößen. Die physikalische Interpretation erfordert einen
	zusätzlichen Schritt: Man muss erkennen, dass $L/R = \tau_L$ die
	Zeitkonstante ist, was nicht unmittelbar evident ist.
	
	\textbf{Mathematische Begründung der Asymmetrie.}
	Die charakteristische Gleichung ist in beiden Fällen identisch:
	\begin{equation}
		s + \frac{R}{L} = 0 \quad \Longrightarrow \quad s = -\frac{R}{L} = -\frac{1}{\tau_L}.
	\end{equation}
	Die Lösung der Differentialgleichung \eqref{eq:rl_standardform} lautet:
	\begin{equation}
		i(t) = \frac{U_0}{R}\left(1 - e^{-t/\tau_L}\right),
		\qquad \tau_L = \frac{L}{R}.
	\end{equation}
	
	\textbf{Energetische Betrachtung.}
	Die im Einschaltvorgang in der Spule gespeicherte Energie berechnet man über:
	\begin{equation}
		W_L(t) = \int_0^t u_L(\tau)\,i(\tau)\,d\tau
		= \int_0^t L\,\dot{i}\,i\,d\tau
		= \frac{1}{2}\,L\,i(t)^2.
	\end{equation}
	
	\begin{figure}[h]
		\centering
		\begin{subfigure}[b]{0.48\textwidth}
			\centering
			\begin{minipage}[c][6.5cm][c]{\linewidth}
				\centering
				\begin{circuitikz}[scale=1.0, transform shape]
					\draw (0,0) to[battery1, l=$U_0$, invert] (0,3)
					to[R, l=$R$] (3,3)
					to[L, l=$L$, v=$u_L$] (3,0)
					to[short] (0,0);
				\end{circuitikz}
			\end{minipage}
			\caption{RL-Einschaltschaltung}
		\end{subfigure}
		\hfill
		\begin{subfigure}[b]{0.48\textwidth}
			\centering
			\begin{minipage}[c][6.5cm][c]{\linewidth}
				\centering
				\begin{tikzpicture}
					\begin{axis}[
						ylabel={Normierte Größe},
						xmin=0, xmax=5,
						ymin=0, ymax=1.1,
						width=7cm, height=5cm,
						grid=major,
						legend pos=north east,
						]
						\addplot[red, thick, domain=0:5, samples=100] {1 - exp(-x)};
						\addlegendentry{$i_L(t)\cdot R/U_0$}
						\addplot[blue, thick, domain=0:5, samples=100] {exp(-x)};
						\addlegendentry{$u_L(t)/U_0$}
					\end{axis}
				\end{tikzpicture}
			\end{minipage}
			\caption{Schwerpunkt: $i_L$ als natürliche Variable}
		\end{subfigure}
		\caption{RL-Einschaltschaltung. Die Spule als Schwerpunkt: $L$ ist direkt der Leitkoeffizient. Die Zeitkonstante $\tau_L = L/R$ ist unmittelbar ablesbar.}
		\label{fig:rl_schwerpunkt}
	\end{figure}
	
	\subsection{Der LC-Parallelschwingkreis: Die Spule als Schwerpunkt der Dynamik}
	\label{sec:schwerpunkt_lc}
	
	\subsubsection*{Die Schaltung}
	
	Im verlustfreien LC-Parallelschwingkreis (vgl.\ Abschnitt~10) sind Spule und
	Kondensator parallel geschaltet. An beiden liegt dieselbe Spannung $u$. Die
	Knotenregel liefert:
	\begin{equation}
		i_C + i_L = 0 \quad \Longrightarrow \quad
		C\,\frac{du}{dt} + i_L = 0.
		\label{eq:lc_knoten}
	\end{equation}
	Da $u = L\,\frac{di_L}{dt}$, ergibt einmalige Differentiation von
	\eqref{eq:lc_knoten}:
	\begin{equation}
		C\,L\,\ddot{i}_L + i_L = 0.
		\label{eq:lc_uc_form}
	\end{equation}
	
	\subsubsection*{Warum die Spule der Schwerpunkt ist}
	
	Die beiden möglichen Formulierungen lauten:
	\begin{align}
		\text{Ãœber } i_L{:}\quad &
		L\,\ddot{i}_L + \frac{1}{C}\,i_L = 0, \label{eq:lc_spule_form} \\
		\text{Ãœber } u_C{:}\quad &
		C\,\ddot{u}_C + \frac{1}{L}\,u_C = 0. \label{eq:lc_kond_form}
	\end{align}
	Division durch den Leitkoeffizienten ergibt in beiden Fällen:
	\begin{equation}
		\ddot{x} + \underbrace{\frac{1}{LC}}_{\omega_0^2}\,x = 0,
		\qquad \omega_0 = \frac{1}{\sqrt{LC}}.
		\label{eq:lc_standardform}
	\end{equation}
	
	\textbf{Mathematische Begründung über den Energiefluss.}
	Die Gesamtenergie des LC-Kreises ist:
	\begin{equation}
		W_{\mathrm{ges}} = W_L + W_C
		= \frac{1}{2}\,L\,i_L^2 + \frac{1}{2}\,C\,u_C^2 = \mathrm{const}.
		\label{eq:lc_energie}
	\end{equation}
	Die Zeitableitung der Gesamtenergie muss verschwinden:
	\begin{equation}
		\dot{W}_{\mathrm{ges}}
		= L\,i_L\,\dot{i}_L + C\,u_C\,\dot{u}_C
		= i_L\,u_L + u_C\,i_C
		= i_L\,u + u\,(-i_L) = 0. \checkmark
	\end{equation}
	
	\textbf{Lösung und Phasenverschiebung.}
	\begin{equation}
		i_L(t) = \hat{I}\,\sin(\omega_0 t), \qquad
		u_C(t) = \hat{U}\,\cos(\omega_0 t),
		\qquad \hat{U} = \hat{I}\sqrt{\frac{L}{C}}.
		\label{eq:lc_loesung}
	\end{equation}
	
	\subsection{Das RL-Glied als Tiefpassfilter: Die Spule als Schwerpunkt der Frequenzselektivität}
	\label{sec:schwerpunkt_rl_tiefpass}
	
	\subsubsection*{Die Schaltung}
	
	Ein Widerstand $R$ in Reihe mit einer Spule $L$ bildet -- wenn die Ausgangsspannung
	über $R$ abgegriffen wird -- einen Tiefpassfilter. Die Masche liefert:
	\begin{equation}
		L\,\frac{di}{dt} + R\,i = u_{in}(t),
		\qquad u_{out} = R\,i.
		\label{eq:rl_tp_masche}
	\end{equation}
	
	\subsubsection*{Übertragungsfunktion -- abgeleitet über die Spule}
	
	Im Frequenzbereich ($j\omega$):
	\begin{equation}
		H_{TP}(j\omega) = \frac{U_{out}}{U_{in}}
		= \frac{R}{R + j\omega L}
		= \frac{1}{1 + j\omega\,\tau_L},
		\qquad \tau_L = \frac{L}{R}.
		\label{eq:rl_tp_uef}
	\end{equation}
	Der Betrag lautet:
	\begin{equation}
		|H_{TP}(j\omega)| = \frac{1}{\sqrt{1 + (\omega\tau_L)^2}}.
		\label{eq:rl_tp_betrag}
	\end{equation}
	
	\textbf{Mathematische Begründung der Frequenzselektivität.}
	\begin{equation}
		|H_{TP}|^2 = \frac{R^2}{R^2 + (\omega L)^2}.
	\end{equation}
	
	\begin{figure}[h]
		\centering
		\begin{subfigure}[b]{0.48\textwidth}
			\centering
			\begin{minipage}[c][6.5cm][c]{\linewidth}
				\centering
				\begin{circuitikz}[scale=1.0, transform shape]
					\draw (0,0) to[short, o-] (0.5,0)
					to[L, l=$L$] (3,0)
					to[R, l=$R$, v=$u_{out}$] (3,-2)
					to[short] (0.5,-2)
					to[short, -o] (0,-2);
					\draw (0,0) to[open, v=$u_{in}$] (0,-2);
					\draw (3,0) to[short, -o] (4.2,0) node[right]{$u_{out}$};
					\draw (3,-2) to[short, -o] (4.2,-2) node[right]{GND};
				\end{circuitikz}
			\end{minipage}
			\caption{RL-Tiefpassfilter}
		\end{subfigure}
		\hfill
		\begin{subfigure}[b]{0.48\textwidth}
			\centering
			\begin{minipage}[c][6.5cm][c]{\linewidth}
				\centering
				\begin{tikzpicture}
					\begin{axis}[
						ylabel={$|H_{TP}|$},
						xmin=0.01, xmax=100,
						ymin=0, ymax=1.1,
						width=7cm, height=5cm,
						xmode=log,
						grid=major,
						legend pos=south west,
						]
						\addplot[blue, thick, domain=0.01:100, samples=200]
						{1/sqrt(1+x^2)};
						\addlegendentry{$|H_{TP}|$}
						\addplot[red, dashed, domain=0.01:1] {1};
						\addplot[red, dashed, domain=1:100] {1/x};
						\addlegendentry{Asymptoten}
						\addplot[only marks, mark=*, red] coordinates {(1,0.707)};
					\end{axis}
				\end{tikzpicture}
			\end{minipage}
			\caption{Bode-Kennlinie des RL-Tiefpasses}
		\end{subfigure}
		\caption{RL-Tiefpassfilter: Grenzfrequenz $\omega_g = R/L$ direkt aus der Spule ablesbar.}
		\label{fig:rl_tiefpass}
	\end{figure}
	
	\subsection{Der ideale Transformator: Die Spule als Schwerpunkt der Kopplung}
	\label{sec:schwerpunkt_trafo}
	
	Die Grundgleichungen des Transformators lauten:
	\begin{align}
		u_1 &= L_1\,\frac{di_1}{dt} + M\,\frac{di_2}{dt}, \label{eq:trafo_prim} \\
		u_2 &= M\,\frac{di_1}{dt} + L_2\,\frac{di_2}{dt}, \label{eq:trafo_sek}
	\end{align}
	wobei $M = \sqrt{L_1 L_2}$ die Gegeninduktivität ist.
	
	Für den idealen Transformator ($k = 1$, verlustfrei):
	\begin{equation}
		\frac{u_1}{u_2} = \frac{N_1}{N_2} = \ddot{u}, \qquad
		\frac{i_1}{i_2} = \frac{N_2}{N_1} = \frac{1}{\ddot{u}},
		\label{eq:trafo_ue}
	\end{equation}
	da $u_k = N_k\,\frac{d\Phi}{dt}$ und der magnetische Fluss $\Phi$ für beide
	Wicklungen identisch ist. Daraus folgt sofort die Leistungserhaltung:
	\begin{equation}
		P_2 = u_2\,i_2 = \frac{u_1}{\ddot{u}} \cdot \ddot{u}\,i_1 = u_1\,i_1 = P_1.
	\end{equation}
	
	\subsection{Zusammenfassende mathematische Begründung}
	\label{sec:spule_schwerpunkt_begruendung}
	
	\textbf{Satz (Spule als Schwerpunkt).}
	\textit{In jeder linearen Schaltung, die mindestens eine Spule enthält,
		ist der Spulenstrom $i_L$ diejenige Variable, bezüglich derer die
		Systemdifferentialgleichung die folgende Eigenschaft besitzt: Der Koeffizient
		der höchsten Zeitableitung von $i_L$ ist ausschließlich die Induktivität $L$
		-- eine physikalisch eigenständige, direkt messbare Größe.}
	
	\textbf{Beweis.}
	Die Spule ist durch das Faradaysche Induktionsgesetz definiert:
	\begin{equation}
		u_L(t) = L\,\frac{di_L}{dt} = \frac{d\Phi}{dt},
		\qquad \Phi = L\,i_L.
		\label{eq:faraday}
	\end{equation}
	In jeder Kirchhoffschen Maschengleichung, die eine Spule enthält, erscheint der
	Term $L\,\dot{i}_L$, und $L$ ist der Koeffizient der höchsten Ableitung
	des Spulenstroms. Die Normierung auf Standardform erfordert lediglich die
	Division durch $L$ -- einen einzigen Schritt mit einem einzigen,
	physikalisch klaren Parameter. $\square$
	
	Die vollständige mechanisch-elektrische Analogie ist in Tabelle~\ref{tab:analogie}
	zusammengefasst.
	
	\begin{table}[h]
		\centering
		\caption{Vollständige Analogie: Mechanik $\longleftrightarrow$ Elektrotechnik}
		\begin{tabular}{>{\centering\arraybackslash}p{4cm}
				>{\centering\arraybackslash}p{2cm}
				>{\centering\arraybackslash}p{4cm}}
			\toprule
			\textbf{Mechanik} & \textbf{Analogie} & \textbf{Elektrotechnik} \\
			\midrule
			Masse $m$ & $\longleftrightarrow$ & Induktivität $L$ \\
			Dämpfung $d$ & $\longleftrightarrow$ & Widerstand $R$ \\
			Federkonstante $k$ & $\longleftrightarrow$ & Kehrwert $1/C$ \\
			Geschwindigkeit $\dot{x}$ & $\longleftrightarrow$ & Strom $i$ \\
			Ort $x$ & $\longleftrightarrow$ & Ladung $q$ \\
			Kraft $F$ & $\longleftrightarrow$ & Spannung $U$ \\
			Impuls $p = m\dot{x}$ & $\longleftrightarrow$ & Magnetischer Fluss $\Phi = Li$ \\
			Kinetische Energie $\tfrac{1}{2}m\dot{x}^2$ & $\longleftrightarrow$ &
			Magnet. Energie $\tfrac{1}{2}Li^2$ \\
			\bottomrule
		\end{tabular}
		\label{tab:analogie}
	\end{table}
	
	Die Spule ist das Trägheitselement der Elektrotechnik: Sie erscheint als
	Leitkoeffizient von $\ddot{i}$ und bestimmt, wie \emph{träge} der Strom auf
	äußere Spannungen reagiert. Jede Schaltung dieser Arbeit, die eine Spule
	enthält, bestätigt diesen Grundsatz:
	
	\begin{equation}
		\underbrace{L}_{\mathrm{Trägheit}}\,\ddot{i}
		+ \underbrace{R}_{\mathrm{Reibung}}\,\dot{i}
		+ \underbrace{\frac{1}{C}}_{\mathrm{Federkonstante}}\,i = f(t).
	\end{equation}
	
	\textbf{Energetische Begründung.}
	Die in der Spule gespeicherte magnetische Energie
	\begin{equation}
		W_L = \int_0^{i} L\,i'\,di' = \frac{1}{2}\,L\,i^2 \geq 0, \quad \forall\, i, L > 0
	\end{equation}
	ist eine quadratische Funktion der Zustandsgröße $i$ -- analog zur kinetischen
	Energie $\tfrac{1}{2}m\dot{x}^2$. Die Nicht-Negativität ist nicht zufällig,
	sondern erzwungen: durch die lineare Kopplung $\Phi = L\,i$ und die Tatsache,
	dass $L > 0$ eine physikalische Eigenschaft jeder realen Induktivität ist.
	Die Spule ist damit auch energetisch ihr eigener Schwerpunkt.
	
	\textbf{Schlussfolgerung.}
	In allen Schaltungen dieser Arbeit, die eine Spule enthalten -- der RL-Einschalt\-schaltung,
	dem LC-Parallelschwingkreis, dem RLC-Serienschwingkreis und dem Transformator --
	ist die Spule der Schwerpunkt: Der Spulenstrom $i_L$ ist die natürliche
	Zustandsgröße, $L$ ist der direkte Leitkoeffizient, und die physikalischen
	Parameter sind ohne weiteren Normierungsschritt ablesbar.
	
	% ============================================================
	\section{Der Operationsverstärker als Integrator}
	% ============================================================
	
	\subsection{Physikalische Beschreibung und Schaltung}
	
	Ein idealer Operationsverstärker mit Kondensator in der Rückkopplung realisiert mathematische Integration:
	\begin{equation}
		u_{out}(t) = -\frac{1}{RC} \int_0^t u_{in}(\tau)\, d\tau
	\end{equation}
	
	\begin{figure}[h]
		\centering
		\begin{subfigure}[b]{0.48\textwidth}
			\centering
			\begin{minipage}[c][6.5cm][c]{\linewidth}
				\centering
				\begin{circuitikz}[scale=1.0, transform shape]
					\draw (0,1.5) node[op amp, yscale=-1] (opamp) {};
					\draw (opamp.+) to[short] (-1, 2) to[R, l=$R$] (-3,2)
					node[left]{$u_{in}$};
					\draw (opamp.-) to[short] (-1,1) to[short] (-1,3)
					to[C, l=$C$] (1.5,3) to[short] (1.5,1.5)
					to[short] (opamp.out);
					\draw (opamp.-) to[short] (-1,1) to[short] (-1,0.5)
					node[ground]{};
					\draw (opamp.out) to[short, -o] (2.5,1.5) node[right]{$u_{out}$};
				\end{circuitikz}
			\end{minipage}
			\caption{OPV-Integrator}
		\end{subfigure}
		\hfill
		\begin{subfigure}[b]{0.48\textwidth}
			\centering
			\begin{minipage}[c][6.5cm][c]{\linewidth}
				\centering
				\begin{tikzpicture}
					\begin{axis}[
						% xlabel={Zeit $t/\tau$},
						ylabel={Normierte Größe},
						xmin=0, xmax=4,
						ymin=-2.1, ymax=1.1,
						width=7cm, height=5cm,
						grid=major,
						legend pos=south west,
						]
						\addplot[blue, thick, domain=0:4, samples=100] {1};
						\addlegendentry{$u_{in}/U_0 = 1$}
						\addplot[red, thick, domain=0:4, samples=100] {-x};
						\addlegendentry{$u_{out} = -t/\tau$}
					\end{axis}
				\end{tikzpicture}
			\end{minipage}
			\caption{$u$-$i$-Diagramm: Rampe als Integral}
		\end{subfigure}
		\caption{OPV-Integrator: Ein konstantes Eingangssignal erzeugt am Ausgang eine lineare Rampe -- mathematische Integration durch physikalische Schaltung.}
		\label{fig:integrator}
	\end{figure}
	
	\subsection{Physikalisch-mathematische Konsistenz}
	
	Der Kondensator in der Rückkopplung speichert die \"aufgelaufene'' Ladung:
	\begin{equation}
		i_C = C \frac{du_{out}}{dt} = -\frac{u_{in}}{R} \quad \Rightarrow \quad u_{out}(t) = -\frac{1}{RC}\int_0^t u_{in}(\tau)\,d\tau.
	\end{equation}
	
	Die physikalische Schaltung realisiert damit exakt die mathematische Operation der Integration. Das $u$-$i$-Diagramm zeigt: Ein konstantes Eingangssignal (Rechteck) erzeugt eine lineare Rampe am Ausgang -- genau wie $\int 1\,dt = t$. Physik und Mathematik sind hier identisch.
	
	% ============================================================
	\section{Der Operationsverstärker als Differenzierer}
	% ============================================================
	
	\subsection{Physikalische Beschreibung und Schaltung}
	
	Die duale Anordnung -- Kondensator am Eingang, Widerstand in der Rückkopplung -- realisiert mathematische Differentiation:
	\begin{equation}
		u_{out}(t) = -RC \cdot \frac{du_{in}}{dt}
	\end{equation}
	
	\begin{figure}[h]
		\centering
		\begin{subfigure}[b]{0.48\textwidth}
			\centering
			\begin{minipage}[c][6.5cm][c]{\linewidth}
				\centering
				\begin{circuitikz}[scale=1.0, transform shape]
					\draw (0,1.5) node[op amp, yscale=-1] (opamp) {};
					\draw (opamp.+) to[short] (-1, 2) to[C, l=$C$] (-3,2)
					node[left]{$u_{in}$};
					\draw (opamp.-) to[short] (-1,1) to[short] (-1,3)
					to[R, l=$R$] (1.5,3) to[short] (1.5,1.5)
					to[short] (opamp.out);
					\draw (opamp.-) to[short] (-1,1) to[short] (-1,0.5)
					node[ground]{};
					\draw (opamp.out) to[short, -o] (2.5,1.5) node[right]{$u_{out}$};
				\end{circuitikz}
			\end{minipage}
			\caption{OPV-Differenzierer}
		\end{subfigure}
		\hfill
		\begin{subfigure}[b]{0.48\textwidth}
			\centering
			\begin{minipage}[c][6.5cm][c]{\linewidth}
				\centering
				\begin{tikzpicture}
					\begin{axis}[
						% xlabel={Zeit $t/\tau$},
						ylabel={Normierte Größe},
						xmin=0, xmax=6.28,
						ymin=-1.1, ymax=1.1,
						width=7cm, height=5cm,
						grid=major,
						legend pos=north east,
						xtick={0,1.57,3.14,4.71,6.28},
						xticklabels={$0$,$\frac{\pi}{2}$,$\pi$,$\frac{3\pi}{2}$,$2\pi$}
						]
						\addplot[blue, thick, domain=0:6.28, samples=100] {sin(deg(x))};
						\addlegendentry{$u_{in} = \sin(\omega t)$}
						\addplot[red, thick, domain=0:6.28, samples=100] {-cos(deg(x))};
						\addlegendentry{$u_{out} = -\cos(\omega t)$}
					\end{axis}
				\end{tikzpicture}
			\end{minipage}
			\caption{Sinus wird zu (negativem) Kosinus}
		\end{subfigure}
		\caption{OPV-Differenzierer: Die Ableitung des Sinus ist der Kosinus. Die physikalische Schaltung realisiert exakt $\frac{d}{dt}\sin(\omega t) = \omega\cos(\omega t)$.}
		\label{fig:differenzierer}
	\end{figure}
	
	\subsection{Physikalisch-mathematische Konsistenz}
	
	Das $u$-$i$-Diagramm demonstriert eindrucksvoll: Die Ableitung eines Sinussignals ist mathematisch der Kosinus, und genau dieses Verhalten zeigt die physikalische Schaltung. Die Phase verschiebt sich um $90^\circ$ -- exakt wie die Mathematik vorhersagt: $\frac{d}{dt}\sin(\omega t) = \omega\cos(\omega t)$.
	
	% ============================================================
	\section{Der Tiefpassfilter: Frequenzabhängige Dämpfung}
	% ============================================================
	
	\subsection{Physikalische Beschreibung und Schaltung}
	
	Ein RC-Tiefpassfilter dämpft hohe Frequenzen und lässt niedrige passieren. Die Übertragungsfunktion lautet:
	\begin{equation}
		H(j\omega) = \frac{1}{1 + j\omega RC} = \frac{1}{1 + j\omega\tau}
	\end{equation}
	
	\begin{figure}[h]
		\centering
		\begin{subfigure}[b]{0.48\textwidth}
			\centering
			\begin{minipage}[c][6.5cm][c]{\linewidth}
				\centering
				\begin{circuitikz}[scale=1.0, transform shape]
					\draw (0,2) to[short, o-] (0.5,2)
					to[R, l=$R$] (2.5,2)
					to[short] (3,2)
					to[C, l=$C$, v=$u_{out}$] (3,0)
					to[short] (0,0)
					to[short, -o] (0,0);
					\draw (0,2) to[open, v=$u_{in}$] (0,0);
					\draw (3,2) to[short, -o] (4,2) node[right]{$u_{out}$};
					\draw (3,0) to[short, -o] (4,0) node[right]{GND};
					% \node[below] at (1.5,-0.4) {RC-Tiefpass};
				\end{circuitikz}
			\end{minipage}
			\caption{RC-Tiefpassschaltung}
		\end{subfigure}
		\hfill
		\begin{subfigure}[b]{0.48\textwidth}
			\centering
			\begin{minipage}[c][6.5cm][c]{\linewidth}
				\centering
				\begin{tikzpicture}
					\begin{axis}[
						% xlabel={Frequenz $f / f_g$ (normiert)},
						ylabel={Betrag $|H|$},
						xmin=0.01, xmax=100,
						ymin=0, ymax=1.1,
						xmode=log,
						width=7cm, height=5cm,
						grid=major,
						legend pos=south west,
						]
						\addplot[blue, thick, domain=0.01:100, samples=200]
						{1/sqrt(1 + x^2)};
						\addlegendentry{$|H| = 1/\sqrt{1+(\omega\tau)^2}$}
						\addplot[red, dashed, domain=0.01:100]
						{(x < 1) ? 1 : 1/x};
						\addlegendentry{Asymptoten}
						\addplot[only marks, mark=*, red] coordinates {(1,0.707)};
						\node at (axis cs:1.5,0.8) {$-3\,\text{dB}$};
					\end{axis}
				\end{tikzpicture}
			\end{minipage}
			\caption{Bode-Betragskennlinie des Tiefpasses}
		\end{subfigure}
		\caption{RC-Tiefpassfilter: Bei der Grenzfrequenz $f_g = 1/(2\pi RC)$ beträgt die Dämpfung exakt $1/\sqrt{2} \approx 0{,}707$ ($-3\,\text{dB}$). Physik und Mathematik stimmen überein.}
		\label{fig:tiefpass}
	\end{figure}
	
	\subsection{Physikalisch-mathematische Konsistenz}
	
	Der Betrag der Ãœbertragungsfunktion $|H(j\omega)| = 1/\sqrt{1+(\omega\tau)^2}$ liegt stets im Intervall $(0,1]$:
	\begin{equation}
		0 < |H(j\omega)| \leq 1, \quad \forall \omega \geq 0.
	\end{equation}
	
	Ein passiver Tiefpassfilter kann das Signal niemals verstärken -- er kann es nur dämpfen. Die Mathematik garantiert dies durch die Ungleichung $\sqrt{1+(\omega\tau)^2} \geq 1$. Physikalisch bedeutet dies: Passive RC-Netzwerke können keine Energie erzeugen.
	
	% ============================================================
	\section{Der Hochpassfilter}
	% ============================================================
	
	\subsection{Physikalische Beschreibung und Schaltung}
	
	Der RC-Hochpassfilter ist die duale Struktur zum Tiefpass -- der Kondensator und Widerstand werden vertauscht:
	\begin{equation}
		H_{HP}(j\omega) = \frac{j\omega\tau}{1 + j\omega\tau}, \qquad |H_{HP}| = \frac{\omega\tau}{\sqrt{1+(\omega\tau)^2}}
	\end{equation}
	
	\begin{figure}[h]
		\centering
		\begin{subfigure}[b]{0.48\textwidth}
			\centering
			\begin{minipage}[c][6.5cm][c]{\linewidth}
				\centering
				\begin{circuitikz}[scale=1.0, transform shape]
					\draw (0,2) to[short, o-] (0.5,2)
					to[C, l=$C$] (2.5,2)
					to[short] (3,2)
					to[R, l=$R$, v=$u_{out}$] (3,0)
					to[short] (0,0)
					to[short, -o] (0,0);
					\draw (0,2) to[open, v=$u_{in}$] (0,0);
					\draw (3,2) to[short, -o] (4,2) node[right]{$u_{out}$};
					\draw (3,0) to[short, -o] (4,0) node[right]{GND};
					% \node[below] at (1.5,-0.4) {RC-Hochpass};
				\end{circuitikz}
			\end{minipage}
			\caption{RC-Hochpassschaltung}
		\end{subfigure}
		\hfill
		\begin{subfigure}[b]{0.48\textwidth}
			\centering
			\begin{minipage}[c][6.5cm][c]{\linewidth}
				\centering
				\begin{tikzpicture}
					\begin{axis}[
						% xlabel={Frequenz $f/f_g$ (normiert)},
						ylabel={Betrag $|H_{HP}|$},
						xmin=0.01, xmax=100,
						ymin=0, ymax=1.1,
						xmode=log,
						width=7cm, height=5cm,
						grid=major,
						legend pos=south east,
						]
						\addplot[blue, thick, domain=0.01:100, samples=200]
						{x/sqrt(1 + x^2)};
						\addlegendentry{$|H_{HP}|$}
						\addplot[red, dashed, domain=0.01:100]
						{(x > 1) ? 1 : x};
						\addlegendentry{Asymptoten}
						\addplot[only marks, mark=*, red] coordinates {(1,0.707)};
					\end{axis}
				\end{tikzpicture}
			\end{minipage}
			\caption{Bode-Betragskennlinie des Hochpasses}
		\end{subfigure}
		\caption{RC-Hochpassfilter: Tiefe Frequenzen werden gedämpft, hohe durchgelassen. Die Betragskennlinie ist die mathematische Duale zum Tiefpass: $|H_{HP}|^2 + |H_{TP}|^2 = 1$.}
		\label{fig:hochpass}
	\end{figure}
	
	\subsection{Physikalisch-mathematische Konsistenz}
	
	Bemerkenswert ist die mathematische Ergänzung von Tief- und Hochpass:
	\begin{equation}
		|H_{TP}(j\omega)|^2 + |H_{HP}(j\omega)|^2 = \frac{1}{1+(\omega\tau)^2} + \frac{(\omega\tau)^2}{1+(\omega\tau)^2} = 1.
	\end{equation}
	
	Physikalisch bedeutet dies: Was der Tiefpass durchlässt, sperrt der Hochpass, und umgekehrt. Die Summe der Leistungen am Tief- und Hochpassausgang ergibt stets die Eingangsleistung -- Energieerhaltung. Die Mathematik drückt dies durch die Identität $a + b = 1$ aus, wobei $a = 1/(1+x^2)$ und $b = x^2/(1+x^2)$.
	
	% ============================================================
	\section{Die Z\"{e}nerdiode als Spannungsbegrenzer}
	% ============================================================
	
	\subsection{Physikalische Beschreibung und Schaltung}
	
	Eine Z\"{e}nerdiode begrenzt die Ausgangsspannung auf den Z\"{e}nerdurchbruchswert $U_Z$:
	\begin{equation}
		u_{out} = \min(u_{in},\; U_Z)
	\end{equation}
	
	\begin{figure}[h]
		\centering
		\begin{subfigure}[b]{0.48\textwidth}
			\centering
			\begin{minipage}[c][6.5cm][c]{\linewidth}
				\centering
				\begin{circuitikz}[scale=1.0, transform shape]
					\draw (0,0) to[short, o-] (1,0)
					to[R, l=$R$] (3,0)
					to[short] (4,0)
					to[short, -o] (5,0) node[right]{$u_{out}$};
					\draw (4,0) to[zD, l=$U_Z$] (4,-2)
					to[short] (1,-2)
					to[short] (0,-2)
					to[short, -o] (0,-2);
					\draw (0,0) to[open, v=$u_{in}$] (0,-2);
				\end{circuitikz}
			\end{minipage}
			\caption{Z\"{e}nerdioden-Begrenzer}
		\end{subfigure}
		\hfill
		\begin{subfigure}[b]{0.48\textwidth}
			\centering
			\begin{minipage}[c][6.5cm][c]{\linewidth}
				\centering
				\begin{tikzpicture}
					\begin{axis}[
						% xlabel={Eingangsspannung $u_{in}$},
						ylabel={Ausgangsspannung $u_{out}$},
						xmin=-1, xmax=5,
						ymin=-1, ymax=3.5,
						width=7cm, height=5cm,
						grid=major,
						legend pos=north west,
						]
						\addplot[blue, thick, domain=-1:2.5, samples=50] {x};
						\addlegendentry{$u_{out} = u_{in}$ (linear)}
						\addplot[red, thick, domain=2.5:5, samples=50] {2.5};
						\addlegendentry{$u_{out} = U_Z$ (Begrenzung)}
						\addplot[dashed, gray] coordinates {(2.5,-1)(2.5,3.5)};
						\node at (axis cs:2.7,0.3) {$U_Z$};
					\end{axis}
				\end{tikzpicture}
			\end{minipage}
			\caption{Kennlinie: Lineare Zone und Begrenzung}
		\end{subfigure}
		\caption{Z\"{e}nerdioden-Spannungsbegrenzer: Unterhalb $U_Z$ linear, bei $U_Z$ Begrenzung. Die Mathematik der Minimum-Funktion beschreibt das physikalische Verhalten exakt.}
		\label{fig:zener}
	\end{figure}
	
	\subsection{Physikalisch-mathematische Konsistenz}
	
	Die Z\"{e}nerdiode implementiert die Minimum-Funktion $u_{out} = \min(u_{in}, U_Z)$. Physikalisch: Im Durchbruchbereich fließt beliebig viel Strom durch die Z\"{e}nerdiode, um die Spannung auf $U_Z$ zu begrenzen. Mathematisch entspricht dies der Klemmung einer Funktion nach oben. Die Nicht-Linearität der physikalischen Schaltung spiegelt sich exakt in der Nicht-Linearität der Minimum-Funktion wider.
	
	% ============================================================
	\section{Berechenbarkeit und Beschreibbarkeit}
	\label{sec:berechenbarkeit}
	% ============================================================

	\subsection{Die Berechenbarkeit}

	Die bisherigen Kapitel dieser Arbeit haben gezeigt, dass physikalische Strukturen dann mathematisch vollständig beschreibbar sind, wenn die zugrundeliegenden Kopplungsrelationen zwischen den beteiligten Größen bekannt und in eine geschlossene mathematische Form gefasst werden können. Diode, Kondensator, Spule, Schwingkreis -- all diese Systeme besitzen eine innere Ordnung, die sich präzise in Gleichungen ausdrücken lässt. Doch damit stellt sich unweigerlich eine tiefere Frage: Ist dies für \textit{jede} physikalische Erscheinung möglich? Ist jede Funktion, die ein Naturphänomen beschreiben soll, tatsächlich beschreibbar -- und wenn ja, ist sie dann auch berechenbar?

	Diese Frage ist nicht trivial. Sie berührt die Grenzen der Mathematik selbst, die Grenzen des Modellierens und letztlich die Frage, was es bedeutet, ein Naturphänomen zu \textit{verstehen}.

	\subsection{Beschreibbarkeit und Berechenbarkeit}

	In dieser Arbeit gilt der folgende grundlegende Zusammenhang als Leitprinzip: \textit{Was beschreibbar ist, ist berechenbar.}

	Dieser Satz klingt zunächst selbstverständlich, ist aber von erheblicher Tragweite. Eine Funktion $f: X \to Y$ ist \textit{beschreibbar}, wenn es eine endliche, eindeutige, regelgeleitete Vorschrift gibt, die jedem Element $x \in X$ ein wohlbestimmtes Element $f(x) \in Y$ zuordnet -- sei es durch eine algebraische Formel, eine Differentialgleichung, ein Grenzwertverfahren oder ein Algorithmus. Eine Funktion ist \textit{berechenbar}, wenn eine Turingmaschine (oder ein äquivalentes Berechnungsmodell) existiert, die für jede Eingabe $x$ nach endlich vielen Schritten die Ausgabe $f(x)$ liefert.

	Der entscheidende Punkt ist: Jede Funktion, die durch eine endliche, explizite mathematische Beschreibung gegeben ist -- eine Formel, eine Gleichung, ein Algorithmus -- definiert zugleich ein Berechnungsverfahren. Die Beschreibung \textit{ist} das Berechnungsverfahren. Beschreibbarkeit impliziert daher Berechenbarkeit. Umgekehrt gilt: Eine Funktion, die nicht beschreibbar ist -- für die also keine endliche, regelgeleitete Zuordnungsvorschrift existiert -- kann auch nicht berechenbar sein. Denn ohne Beschreibung gibt es kein Verfahren, das eine Maschine ausführen könnte.

	Formaler gesagt: Die Menge der beschreibbaren Funktionen und die Menge der berechenbaren Funktionen fallen zusammen. Was nicht in einer endlichen symbolischen Sprache ausgedrückt werden kann, existiert als Berechnungsvorschrift nicht.

	\subsection{Was ist nicht immer beschreibbar?}

	Die eigentlich interessante und herausfordernde Frage ist nicht, ob es nicht-berechenbare Funktionen \textit{im mathematischen Sinne} gibt -- das ist durch den Satz von Cantor und die Ergebnisse von Turing, Church und Gödel seit dem 20.~Jahrhundert bekannt. Die tiefere Frage lautet: Welche Funktionen, die \textit{physikalische Phänomene} beschreiben sollen, sind \textit{nicht immer} beschreibbar?

	Die Antwort liegt in der Natur bestimmter physikalischer Phänomene selbst. Betrachten wir Prozesse, die nicht kontinuierlich verlaufen, sondern sprunghaft, chaotisch, diskontinuierlich oder prinzipiell nicht vollständig vorhersagbar sind:

	\begin{itemize}
		\item \textbf{Explosionen und Detonationen:} Eine Explosion ist ein physikalischer Prozess, bei dem sich Energie in extrem kurzer Zeit ungeordnet und hochgradig nichtlinear entlädt. Die Druckwelle, die Temperaturverteilung, die Fragmentierung -- all diese Größen hängen von einer Vielzahl von Parametern ab, die in ihrer Gesamtheit weder vollständig messbar noch vollständig modellierbar sind. Eine Funktion, die den genauen zeitlichen und räumlichen Verlauf einer Explosion für alle möglichen Anfangsbedingungen beschreiben soll, ist \textit{nicht in geschlossener Form angebbar}. Es existiert keine Formel $f(t, \mathbf{x}, \text{Anfangsbedingungen})$, die diesen Verlauf für alle denkbaren Szenarien exakt und vollständig beschreibt.
		\item \textbf{Vulkanausbrüche und geophysikalische Diskontinuitäten:} Ein Vulkanausbruch ist ein paradigmatisches Beispiel für ein physikalisches Phänomen, das diskontinuierlich verläuft. Der Übergang vom ruhenden Vulkan zum aktiven Ausbruch ist kein stetiger, differenzierbarer Prozess. Druckaufbau, Magmafließen, seismische Aktivität -- diese Größen zeigen Sprünge, Instabilitäten und Bifurkationen, die durch partielle Differentialgleichungen nur näherungsweise und unter starken vereinfachenden Annahmen erfasst werden können. Eine \textit{universelle} mathematische Funktion, die den vollständigen Verlauf eines Vulkanausbruchs für beliebige Bedingungen beschreibt, existiert nicht.
		\item \textbf{Turbulenz:} Die turbulente Strömung ist ein weiteres klassisches Beispiel. Obwohl die Navier-Stokes-Gleichungen die physikalischen Grundgesetze der Strömungsmechanik vollständig beschreiben, ist die analytische Lösung für turbulente Regime nicht bekannt -- und es ist ein offenes Millennium-Problem der Mathematik, ob für dreidimensionale turbulente Strömungen überhaupt reguläre, eindeutige Lösungen existieren. Die Funktion, die das Geschwindigkeitsfeld einer turbulenten Strömung vollständig beschreibt, ist damit nicht geschlossen angebbar.
		\item \textbf{Chaotische Systeme mit sensitiver Anfangsbedingungsabhängigkeit:} Ein deterministisch chaotisches System wie das Lorenz-System $\dot{x} = \sigma(y-x)$, $\dot{y} = x(\rho - z) - y$, $\dot{z} = xy - \beta z$ ist \textit{im Prinzip} beschreibbar -- es gibt eine Differentialgleichung. Aber die Funktion, die den Zustand des Systems zu einem späteren Zeitpunkt $t$ aus dem Anfangszustand berechnet, ist nicht mehr praktisch beschreibbar, weil beliebig kleine Änderungen der Anfangsbedingungen exponentiell verstärkt werden (positive Lyapunov-Exponenten). Die theoretische Beschreibbarkeit kollabiert in der praktischen Unberechenbarkeit.
	\end{itemize}

	Das gemeinsame Merkmal dieser Phänomene ist: Sie sind \textit{nicht vollständig kontinuierlich beschreibbar}. Ihr physikalischer Verlauf enthält Diskontinuitäten, Sprünge, sensitive Abhängigkeiten oder prinzipielle Informationsverluste, die eine geschlossene analytische Beschreibung unmöglich machen oder zumindest grundlegend einschränken.

	\subsection{Nicht beschreibbar ist nicht berechenbar}

	Aus dem in Abschnitt~\ref{sec:berechenbarkeit} formulierten Prinzip -- Was beschreibbar ist, ist berechenbar; was nicht beschreibbar ist, ist nicht berechenbar -- folgt direkt: \textit{Wenn eine Funktion nicht beschreibbar ist, dann ist sie auch nicht berechenbar.}  Und ebenso gilt die Umkehrung: \textit{Wenn eine Funktion nicht berechenbar ist, dann ist sie auch nicht mehr beschreibbar, denn alle beschreibbaren Funktionen sind berechenbar.}

	Diese Äquivalenz ist kein Zirkelschluss, sondern ein inhaltliches Prinzip. Eine endliche, explizite Beschreibung einer Funktion enthält stets die Information, wie diese Funktion schrittweise ausgewertet werden kann -- das heißt, wie sie berechnet werden kann. Und umgekehrt: Gibt es kein endliches Verfahren zur Berechnung, dann kann es auch keine endliche, explizite Beschreibung geben, denn eine solche würde sofort ein Berechnungsverfahren liefern.

	Die Konsequenz für die physikalische Modellierung ist klar und bedeutsam: Phänomene wie Explosionen, Vulkanausbrüche oder turbulente Strömungen, die nicht vollständig und allgemeingültig in einer geschlossenen mathematischen Funktion beschreibbar sind, können auch nicht vollständig berechnet werden. Numerische Simulationen solcher Phänomene liefern stets nur Näherungen unter spezifischen, idealisierten Bedingungen. Sie sind partiell beschreibbar und partiell berechenbar -- aber nie universal.

	Für den Fall, dass eine Funktion nicht mehr physikalisch oder mathematisch beschreibbar ist, aber berechenbar sein soll, gilt: Es ergibt sich dann eine Folge von Zuständen, für die dann am Ende eine Teilfolge von Zuständen existiert, dass wenn diese Teilfolge während der Berechnung betreten wird, dass es dann einen letzten Zustand in dieser Teilfolge gibt, ab dem die Maschine nicht mehr natürlich weiterrechnen kann.

	Dieser Sachverhalt beschreibt präzise die Situation, die sich an den Grenzen der physikalischen Beschreibbarkeit einstellt: Eine Rechenmaschine verarbeitet eine Folge von wohldefiniierten Zuständen -- solange die zugrundeliegende Funktion beschreibbar ist, schreitet die Berechnung regelgeleitet fort. Sobald jedoch der Bereich betreten wird, in dem die Funktion nicht mehr beschreibbar ist -- weil das Phänomen diskontinuierlich wird, weil Singularitäten auftreten, weil die physikalische Realität keine geschlossene Fortsetzung mehr erlaubt -- existiert innerhalb der durchlaufenen Zustandsfolge eine Teilfolge, die in einen letzten Zustand mündet: einen Zustand, ab dem kein natürlich beschreibbarer nächster Schritt mehr existiert. Die Maschine kann nicht weiter \textit{natürlich} rechnen -- das heißt, nicht weiter nach einer endlichen, durch die Physik gegebenen Beschreibungsregel.

	Dieses Bild ist keine Metapher, sondern eine strukturelle Aussage: Die Berechenbarkeit endet dort, wo die Beschreibbarkeit endet -- und die Beschreibbarkeit endet dort, wo das Naturphänomen aufhört, kontinuierlich und regelgeleitet zu sein.
	
	\subsection{Beschreibbarkeit ist nicht Berechenbarkeit}
	
	Warum sind Beschreibbarkeit und Berechenbarkeit nicht gleich? Die Beschreibbarkeit ist wichtiger, denn von Natur aus versucht der Mensch, die Phänomene in der Natur erst mal intuitiv mit den Gedanken und vielleicht auch mit einem Stift und einem Papier zu beschreiben. Das ist wichtig und wird auch immer so bleiben. Nicht jede Beschreibung muss auch programmatisch berechnet werden. Das wäre übertrieben. Auf der anderen Seite ist erst die Berechenbarkeit der Zugang zu automatisiereten und wirklich aussakekräftigen Ergebnissen von wichtigen Beschreibungen.

	\subsection{Konsequenzen für die physikalische Modellierung}

	Die Einsicht, dass Beschreibbarkeit und Berechenbarkeit äquivalent sind, und dass bestimmte physikalische Phänomene prinzipiell nicht vollständig beschreibbar sind, hat fundamentale Konsequenzen für die wissenschaftliche Praxis:

	\begin{enumerate}
		\item \textbf{Modelle sind stets Näherungen.} Jedes mathematische Modell eines physikalischen Phänomens beschreibt nur den Bereich, in dem das Phänomen durch die gewählte Beschreibungssprache erfasst werden kann. Außerhalb dieses Bereichs versagt das Modell -- nicht wegen mangelnder Rechenleistung, sondern prinzipiell.
		\item \textbf{Die Grenze der Beschreibbarkeit ist eine physikalische Grenze.} Wenn ein Phänomen diskontinuierlich, singulär oder chaotisch wird, ist dies keine Frage der Rechenkapazität einer Maschine. Es ist eine Frage der physikalischen Natur des Phänomens selbst. Die Mathematik kann nur beschreiben, was die Physik ihr anbietet.
		\item \textbf{Beschreibbare Phänomene sind ausgezeichnet.} Die in dieser Arbeit betrachteten physikalischen Strukturen -- Widerstand, Kondensator, Spule, Diode, Schwingkreis -- sind gerade deshalb so präzise mathematisch beschreibbar, weil sie in wohldefinierten, stetigen, differenzierbaren Bereichen operieren. Ihre innere Ordnung -- die Kopplungsrelationen -- ist die physikalische Bedingung für ihre Beschreibbarkeit und damit für ihre Berechenbarkeit.
		\item \textbf{Die Äquivalenz als erkenntnistheoretisches Prinzip.} Der Satz \textit{Was beschreibbar ist, ist berechenbar} ist kein technisches Detail der Informatik. Er ist ein erkenntnistheoretisches Prinzip, das aussagt: Erkenntnis entsteht durch Beschreibung, und Beschreibung ist -- wo sie vollständig ist -- Berechnung. Wo die Natur keine vollständige Beschreibung erlaubt, entzieht sie sich auch der vollständigen Erkenntnis.
	\end{enumerate}

	% ============================================================
	\section{Zusammenfassung und allgemeines Prinzip}
	% ============================================================
	
	\subsection{Das allgemeine Muster}
	
	Die betrachteten physikalischen Strukturen weisen ein gemeinsames Muster auf: Immer dann, wenn zwei physikalische Größen, die in einer Produkt- oder Quotientenbeziehung stehen, nicht unabhängig voneinander sind, sondern durch eine physikalische Gesetzmäßigkeit miteinander verknüpft sind, ergibt sich eine quadratische (oder anderweitig nicht-negative) Abhängigkeit in der Energiebeschreibung.
	
	Die Grundstruktur lässt sich wie folgt zusammenfassen:
	\begin{enumerate}
		\item \textbf{Allgemeine Formel:} Eine Formel $E = A \cdot B$, in der $A$ und $B$ scheinbar unabhängig variierbar sind und negative Werte für $E$ ermöglichen würden.
		\item \textbf{Physikalische Kopplung:} Die Struktur erzwingt $A = f(B)$, sodass $A$ und $B$ nicht wirklich unabhängig sind.
		\item \textbf{Mathematische Konsequenz:} Die Einbeziehung der Kopplung ergibt eine nicht-negative Darstellung.
		\item \textbf{Vollständige Konsistenz:} Erst das gemeinsame Bild von Physik und Mathematik ergibt den vollständigen, widersprufsfreien Sinn.
	\end{enumerate}
	
	\subsection{Ãœbersicht aller betrachteten Strukturen}
	
	\begin{table}[h]
		\centering
		\caption{übersicht der physikalischen Strukturen, ihrer Schaltungen und mathematischen Beschreibung}
		\begin{tabular}{>{\raggedright\arraybackslash}p{3cm}>{\raggedright\arraybackslash}p{3cm}>{\raggedright\arraybackslash}p{3cm}>{\raggedright\arraybackslash}p{4cm}}
			\toprule
			\textbf{Struktur} & \textbf{Allg.\ Formel} & \textbf{Kopplungs\-relation} & \textbf{Nicht-negative Form / Eigenschaft} \\
			\midrule
			Diode + Widerstand & $P = U \cdot I$ & $U = R\cdot I$, $I \geq 0$ & $P = R I^2 \geq 0$ \\
			Kondensator (Laden) & $W = Q \cdot U$ & $Q = C \cdot U$ & $W = \tfrac{1}{2}CU^2 \geq 0$ \\
			Spule (Einschalten) & $W = \Phi \cdot I$ & $\Phi = L \cdot I$ & $W = \tfrac{1}{2}LI^2 \geq 0$ \\
			RC-Entladeglied & $P = u \cdot i$ & $u_C = Ri$ (Entladung) & $P_R = Ri^2 \geq 0$ \\
			Halbwellengr. & $u_{out} = u_{in}$ & Diode sperrt $<0$ & $u_{out} = \max(u_{in},0) \geq 0$ \\
			Brückengleichr. & $u_{out} = u_{in}$ & 4 Dioden & $u_{out} = |u_{in}| \geq 0$ \\
			Spannungsteiler & $U_{out}/U_{in} = k$ & $R_1, R_2 > 0$ & $0 < k < 1$ \\
			Transformator & $P = U \cdot I$ & $U_2 I_2 = U_1 I_1$ & Leistungserhaltung \\
			LC-Schwingkreis & $W = W_L + W_C$ & $\sin^2 + \cos^2 = 1$ & $W_{ges} = \text{const} \geq 0$ \\
			RLC-Schwingkreis & DGL 2. Ordnung & $\zeta = R/(2\sqrt{L/C})$ & Stabile Lösungen, $W \searrow 0$ \\
			OPV-Integrator & $u_{out} = f(u_{in})$ & $C \dot{u}_{out} = -u_{in}/R$ & $u_{out} = -\tfrac{1}{RC}\int u_{in}\,dt$ \\
			OPV-Differenz. & $u_{out} = f(u_{in})$ & $i_C = C\dot{u}_{in}$ & $u_{out} = -RC\,\dot{u}_{in}$ \\
			Tiefpassfilter & $|H| \leq 1$ & $|H|^2 = 1/(1+(\omega\tau)^2)$ & $0 < |H| \leq 1$ \\
			Hochpassfilter & $|H| \leq 1$ & $|H_{TP}|^2 + |H_{HP}|^2 = 1$ & Komplementarität \\
			Z\"{e}nerdiode & $u_{out} = u_{in}$ & Durchbruch bei $U_Z$ & $u_{out} = \min(u_{in}, U_Z)$ \\
			\bottomrule
		\end{tabular}
		\label{tab:uebersicht}
	\end{table}
	
	\subsection{Schlussbetrachtung}
	
	Wichtig in dieser Arbeit sind all die physikalischen Strukturen, in denen sich Physik und Mathematik, wenn sie beide betrachtet werden, den vollständigen Sinn für die physikalische Struktur mit ihrer Beschreibung ergeben.
	
	Die betrachteten Beispiele zeigen eindrucksvoll, dass Physik und Mathematik keine getrennten Welten sind. Die Mathematik liefert das formale Werkzeug -- Gleichungen, Funktionen, Ableitungen, Übertragungsfunktionen -- und die Physik liefert die Bedeutung, die Randbedingungen und die Interpretation. Erst wenn beide Perspektiven zusammengebracht werden, ergibt sich das vollständige, widerspruüchsfreie Bild.
	
	Zu jedem System wurde in dieser Arbeit nicht nur die mathematische Beschreibung gegeben, sondern auch die zugehörige elektronische Schaltung (gezeichnet mit \texttt{circuitikz}) sowie das $u$-$i$-Diagramm, das die zeitlichen Verläufe von Spannung und Strom zeigt. Diese dreifache Darstellung -- Schaltung, Diagram, Formel -- macht die Einheit von Physik und Mathematik unmittelbar sichtbar und anschaulich.
	
	Diese wechselseitige Ergänzung von Physik und Mathematik ist eines der schönsten und tiefgründigsten Prinzipien der Naturwissenschaften.

	% ============================================================
	\subsection{Ausblick}
	% ============================================================

	Die in dieser Arbeit entwickelten Grundprinzipien -- insbesondere das Zusammenspiel von physikalischer Kopplungsrelation und mathematischer Nicht-Negativität -- sind keineswegs auf die betrachteten Schaltungen beschränkt. Im Gegenteil: Sie bilden den Ausgangspunkt für eine Vielzahl tiefgreifender Erweiterungen, die in unterschiedlichste Gebiete der Physik, Mathematik und Technik hineinreichen. Der folgende Ausblick entfaltet systematisch, wohin die hier gewonnenen Einsichten führen können.

	% ---
	\subsubsection*{1. \quad Nichtlineare Systeme und verallgemeinerte Kopplungsrelationen}

	Alle in dieser Arbeit betrachteten Schaltungen sind \textit{linear}: Die Kopplungsrelation zwischen zwei physikalischen Größen ist stets proportional -- $U = RI$, $Q = CU$, $\Phi = LI$. Die eigentliche Natur vieler physikalischer Bauelemente ist jedoch nichtlinear.

	Das wichtigste Beispiel ist die Halbleiterdiode. Ihre Strom-Spannungs-Kennlinie folgt nicht dem Ohmschen Gesetz, sondern der \textit{Shockley-Gleichung}:
	\begin{equation}
		I(U) = I_S \left(\mathrm{e}^{U/(n U_T)} - 1\right),
	\end{equation}
	wobei $I_S$ der Sättigungsstrom, $n$ der Idealitätsfaktor und $U_T = k_B T / q$ die Temperaturspannung ist. Trotz dieser Nichtlinearität bleibt das Grundprinzip erhalten: Die am Bauelement umgesetzte Leistung
	\begin{equation}
		P(U) = U \cdot I_S \left(\mathrm{e}^{U/(n U_T)} - 1\right)
	\end{equation}
	ist für $U > 0$ positiv und für $U < 0$ negativ aber betragsmäßig sehr klein -- das physikalische Verhalten eines passiven Bauelements spiegelt sich weiterhin in der Nichtnegativität der zeitgemittelten Leistungsaufnahme. Die Herausforderung bei nichtlinearen Systemen liegt darin, eine verallgemeinerte \textit{Energiefunktion} (eine sogenannte \textit{Lyapunov-Funktion}, vgl.\ Abschnitt 6 dieses Ausblicks) zu finden, die die Rolle der quadratischen Energieausdrücke $\frac{1}{2}CU^2$ und $\frac{1}{2}LI^2$ übernimmt.

	Für nichtlineare Kondensatoren und Spulen verallgemeinern sich die Energieausdrücke zu:
	\begin{equation}
		W_C = \int_0^Q u(q)\,\mathrm{d}q, \qquad W_L = \int_0^\Phi i(\varphi)\,\mathrm{d}\varphi.
	\end{equation}
	Diese Ausdrücke sind genau dann nicht-negativ, wenn die Kennlinien $u(q)$ bzw.\ $i(\varphi)$ streng monoton wachsend durch den Ursprung verlaufen -- eine unmittelbar physikalisch interpretierbare Bedingung für ein passives Verhalten.

	Eine weitere wichtige nichtlineare Struktur ist der \textit{Varaktor} (spannungsabhängiger Kondensator), der in Phasenregelkreisen und HF-Schaltungen Anwendung findet, sowie der \textit{Memristor} -- ein von L.~O.~Chua~\cite{chua1969} theoretisch postuliertes und erst 2008 experimentell realisiertes Bauelement, das die Ladung und den magnetischen Fluss direkt koppelt: $\mathrm{d}\Phi = M(q)\,\mathrm{d}q$, wobei $M(q)$ die \textit{Memristanz} (Gedächtnis-Widerstand) heißt. Der Memristor schließt das von Chua aufgestellte Symmetrieschema der vier grundlegenden Schaltkreisvariablen Spannung, Strom, Ladung und magnetischer Fluss logisch ab. Er verkörpert auf besonders eindrucksvolle Weise das Thema dieser Arbeit: Die Kopplung zweier scheinbar unabhängiger physikalischer Größen erzwingt ein spezifisches, mathematisch beschreibbares Verhalten.

	% ---
	\subsubsection*{2. \quad Mehrtor-Netzwerke und die allgemeine Netzwerktheorie}

	Die in dieser Arbeit betrachteten Strukturen sind fast ausschließlich Eintor-Schaltungen (Zweipole): Es gibt genau ein Klemmpaar, an dem Energie ausgetauscht wird. Die reale Welt der Elektronik und der Physik ist jedoch von \textit{Mehrtor-Netzwerken} geprägt -- Schaltungen mit mehreren Klemmenpaaren, die miteinander interagieren.

	Das einfachste Mehrtor-Netzwerk ist das \textit{Zweitor} (Vierpol), das eine Ein- und eine Ausgangsseite besitzt. Transformator und Operationsverstärker sind bereits Vertreter dieser Klasse. Die allgemeine Beschreibung eines Zweitors erfolgt durch eine $2 \times 2$-Matrix, beispielsweise die \textit{Streumatrix} (S-Matrix) in der Hochfrequenztechnik oder die \textit{Impedanzmatrix} $Z$:
	\begin{equation}
		\begin{pmatrix} U_1 \\ U_2 \end{pmatrix} = \begin{pmatrix} Z_{11} & Z_{12} \\ Z_{21} & Z_{22} \end{pmatrix} \begin{pmatrix} I_1 \\ I_2 \end{pmatrix}.
	\end{equation}
	Für ein passives, verlustfreies Zweitor gilt, dass die Inzidenzmatrix $Z$ \textit{schiefsymmetrisch} sein muss (Reziprozität): $Z_{12} = Z_{21}$. Die Bedingung der Passivität -- dass das Netzwerk keine Netto-Energie erzeugen kann -- lässt sich dann als \textit{positive Semidefinitheit} der Matrix $\mathrm{Re}(Z)$ formulieren:
	\begin{equation}
		\mathbf{I}^* \cdot \mathrm{Re}(Z) \cdot \mathbf{I} \geq 0 \quad \forall\, \mathbf{I} \in \mathbb{C}^2.
	\end{equation}
	Dies ist die direkte Verallgemeinerung der skalaren Bedingung $P = R \cdot I^2 \geq 0$ auf den mehrdimensionalen Fall. Das Konzept der positiv semidefiniten Matrix als matematisches Analogon des Energieprinzips zieht sich durch die gesamte lineare Netzwerktheorie und bildet die Brücke zur Theorie der positiv semidefiniten Matrizen in der linearen Algebra.

	% ---
	\subsubsection*{3. \quad Analogien zwischen elektrischen, mechanischen und thermischen Systemen}

	Eines der folgenreichsten Konzepte der Ingenieursphysik ist das Prinzip der \textit{Systemanalogie}: Schaltungen aus vollständig verschiedenen physikalischen Domänen können durch identische Differentialgleichungen beschrieben werden, wenn die richtigen Größenpaare einander zugeordnet werden.

	Die \textit{Impedanzanalogie} zwischen elektrischen und mechanischen Systemen ist ein klassisches Beispiel:
	\begin{center}
		\begin{tabular}{lll}
			\toprule
			\textbf{Elektrisch} & \textbf{Mechanisch (Kraft-Strom)} & \textbf{Mechanisch (Kraft-Spannung)} \\
			\midrule
			Spannung $u$ & Geschwindigkeit $v$ & Kraft $F$ \\
			Strom $i$ & Kraft $F$ & Geschwindigkeit $v$ \\
			Widerstand $R$ & Dämpfung $b$ & Dämpfung $b$ \\
			Kapazität $C$ & Masse $m$ & Nachgiebigkeit $1/k$ \\
			Induktivität $L$ & Nachgiebigkeit $1/k$ & Masse $m$ \\
			\bottomrule
		\end{tabular}
	\end{center}

	Der LC-Schwingkreis und das Masse-Feder-System sind mathematisch identisch -- beide werden durch eine gewöhnliche Differentialgleichung zweiter Ordnung
	\begin{equation}
		\ddot{x} + \omega_0^2 x = 0, \qquad \omega_0 = \frac{1}{\sqrt{LC}} = \sqrt{\frac{k}{m}}
	\end{equation}
	beschrieben. Die Grundbeobachtung dieser Arbeit über den LC-Kreis -- dass die Gesamtenergie $W_{ges} = \frac{1}{2}LI^2 + \frac{1}{2}CU_C^2 = \mathrm{const.}$ durch $\sin^2 + \cos^2 = 1$ garantiert wird -- gilt wörtlich auch für das Masse-Feder-System: $W_{ges} = \frac{1}{2}mv^2 + \frac{1}{2}kx^2 = \mathrm{const.}$

	Noch weiter reicht die Analogie zu thermischen Systemen. Hier spielen die Rollen:
	\begin{itemize}
		\item Wärmestrom $\dot{Q}$ $\leftrightarrow$ elektrischer Strom $I$
		\item Temperaturdifferenz $\Delta T$ $\leftrightarrow$ Spannung $U$
		\item Wärmewiderstand $R_{th}$ $\leftrightarrow$ elektrischer Widerstand $R$
		\item Wärmekapazität $C_{th}$ $\leftrightarrow$ elektrische Kapazität $C$
	\end{itemize}
	Die instationäre Wärmeleitung in einem Körper verhält sich wie ein RC-Tiefpassfilter: Schnelle Temperaturschwankungen werden gedämpft, langsame passieren nahezu ungehindert. Die in dieser Arbeit für den Tiefpassfilter gefundene Übertragungsfunktion $|H(\omega)| = (1 + (\omega\tau)^2)^{-1/2}$ ist direkt auf das thermische RC-Modell übertragbar, wobei $\tau = R_{th} C_{th}$ die thermische Zeitkonstante ist.

	Diese Analogien illustrieren die universelle Gültigkeit der hier entwickelten Prinzipien: Das Wechselspiel von Kopplungsrelation und Energieform ist kein Spezifikum der Elektrotechnik, sondern ein physikdomänenübergreifendes Strukturprinzip.

	% ---
	\subsubsection*{4. \quad Port-Hamiltonsche Systeme: die allgemeine mathematische Rahmenbeschreibung}

	Der in dieser Arbeit wiederholt beobachtete Zusammenhang zwischen Energieerhaltung, Kopplungsrelationen und der mathematischen Struktur der beschreibenden Gleichungen hat eine präzise und sehr allgemeine mathematische Formulierung: die Theorie der \textit{Port-Hamiltonschen Systeme}.

	Ein Port-Hamiltonsches System ist ein dynamisches System der Form
	\begin{equation}
		\dot{\mathbf{x}} = \bigl(J(\mathbf{x}) - R(\mathbf{x})\bigr)\, \nabla H(\mathbf{x}) + G(\mathbf{x})\, \mathbf{u},
	\end{equation}
	wobei:
	\begin{itemize}
		\item $H(\mathbf{x})$ die \textit{Hamiltonsche Funktion} (Gesamtenergie) des Systems ist,
		\item $J(\mathbf{x})$ eine schiefsymmetrische Matrix ist, die die \textit{verlustfreien} Energieflüsse beschreibt ($J = -J^T$, sodass $\nabla H^T J \nabla H = 0$),
		\item $R(\mathbf{x}) \geq 0$ eine positiv semidefinite Matrix ist, die die \textit{Dissipation} beschreibt,
		\item $G(\mathbf{x})$ die Eingangsmatrix und $\mathbf{u}$ die externe Anregung ist.
	\end{itemize}

	Für ein isoliertes System ($\mathbf{u} = 0$) folgt unmittelbar:
	\begin{equation}
		\dot{H} = \nabla H^T \dot{\mathbf{x}} = \nabla H^T (J - R) \nabla H = \underbrace{\nabla H^T J \nabla H}_{= 0} - \underbrace{\nabla H^T R \nabla H}_{\geq 0} \leq 0.
	\end{equation}
	Die Energie kann nicht zunehmen -- sie bleibt konstant (verlustfreies System) oder nimmt ab (dissipatves System). Dies ist die allgemeinste mathematische Formulierung des Energieprinzips, das diese Arbeit in zahlreichen Einzelbeispielen illustriert.

	Der LC-Schwingkreis dieser Arbeit ist ein Port-Hamiltonsches System mit $H = \frac{1}{2}LI^2 + \frac{1}{2}CU_C^2$, $J = \begin{pmatrix}0 & -1 \\ 1 & 0\end{pmatrix}$ und $R = 0$. Das RLC-Serienglied ist dasselbe System mit dem zusätzlichen Dissipationsterm $R = \begin{pmatrix}R & 0 \\ 0 & 0\end{pmatrix}$. Die Port-Hamiltonsche Theorie macht also transparent, warum der Übergang vom LC- zum RLC-Kreis lediglich das Hinzufügen einer positiv semidefiniten Matrix bedeutet.

	Port-Hamiltonsche Systeme finden Anwendung in der Robotik, der thermodynamischen Systemtheorie, der Plasmatheorie und neuerdings auch in der Modellierung biologischer Netzwerke. Sie bilden die natürliche Sprache für das Thema dieser Arbeit auf dem Stand der modernen mathematischen Physik.

	% ---
	\subsubsection*{5. \quad Variationsprinzipien und Lagrange-Mechanik}

	Die quadratischen Energieformen, die in dieser Arbeit immer wieder auftreten, weisen auf eine tiefere mathematische Struktur hin: die Variationsprinzipien der klassischen Mechanik.

	Das \textit{Hamiltonsche Prinzip der kleinsten Wirkung} besagt, dass die zeitliche Entwicklung eines physikalischen Systems diejenige ist, die das \textit{Wirkungsintegral}
	\begin{equation}
		S = \int_{t_1}^{t_2} \mathcal{L}(q, \dot{q}, t)\,\mathrm{d}t
	\end{equation}
	stationär (in der Regel minimal) macht. Die Lagrange-Funktion $\mathcal{L} = T - V$ ist die Differenz aus kinetischer Energie $T$ und potentieller Energie $V$. Beide Energieformen sind für alle hier betrachteten Systeme nicht-negativ:
	\begin{equation}
		T = \frac{1}{2}LI^2 = \frac{1}{2}m\dot{x}^2 \geq 0, \qquad V = \frac{1}{2}CU_C^{-1}Q^2 = \frac{1}{2}kx^2 \geq 0.
	\end{equation}
	Die Euler-Lagrange-Gleichungen für diese Systeme reproduzieren exakt die in dieser Arbeit hergeleiteten Differentialgleichungen -- für den LC-Schwingkreis:
	\begin{equation}
		\frac{\mathrm{d}}{\mathrm{d}t}\frac{\partial \mathcal{L}}{\partial \dot{Q}} - \frac{\partial \mathcal{L}}{\partial Q} = 0 \implies L\ddot{Q} + \frac{Q}{C} = 0,
	\end{equation}
	was identisch mit $L\ddot{I} + \frac{I}{C} = 0$ (nach Differentiation nach der Zeit) ist.

	Dieser Zusammenhang zeigt, dass die in dieser Arbeit intuitiv entwickelten Kopplungsrelationen ihre tiefste Begründung in einem universellen Variationsprinzip finden. Das Hamiltonsche Prinzip ist dabei nicht spezifisch für Mechanik oder Elektromagnetismus -- es ist ein metatheoretisches Prinzip, das Physik und Mathematik auf der höchsten Abstraktionsebene verbindet.

	Eine Erweiterung führt auf die Theorie der \textit{symmetrischen Systeme} und die Noether-Symmetrien: Jede Symmetrie des Wirkungsintegrals (z.\,B.\ Zeitinvarianz) entspricht nach dem Noetherschen Theorem einer Erhaltungsgröße (z.\,B.\ Energie). Die Energieerhaltung im LC-Schwingkreis ist daher keine Zufälligkeit, sondern Konsequenz der Zeitinvarianz der zugrundeliegenden physikalischen Gesetze.

	% ---
	\subsubsection*{6. \quad Stabilität, Lyapunov-Theorie und Systemanalyse}

	Der Begriff der Energiefunktion, der in dieser Arbeit implizit stets präsent war (Energie eines Kondensators, Energie einer Spule, Gesamtenergie eines Schwingkreises), hat in der modernen Kontrolltheorie eine präzise mathematische Entsprechung: die \textit{Lyapunov-Funktion}.

	Ein autonomes dynamisches System $\dot{\mathbf{x}} = f(\mathbf{x})$ mit Gleichgewichtspunkt $\mathbf{x}^* = 0$ ist asymptotisch stabil, wenn eine stetig differenzierbare Funktion $V(\mathbf{x}) > 0$ mit $V(0) = 0$ existiert, sodass
	\begin{equation}
		\dot{V}(\mathbf{x}) = \nabla V(\mathbf{x})^T f(\mathbf{x}) < 0 \quad \forall\, \mathbf{x} \neq 0.
	\end{equation}
	Diese Bedingung ist für alle in dieser Arbeit betrachteten gedämpften Systeme erfüllt, wenn man $V = H$ (die Gesamtenergie) als Lyapunov-Funktion wählt: Die Energie nimmt monoton ab, bis das System zur Ruhe kommt.

	Besonders aufschlussreich ist das RLC-Serienglied: Die Gesamtenergie $H = \frac{1}{2}LI^2 + \frac{1}{2}CU_C^2$ ist positiv definit, und ihre Zeitableitung
	\begin{equation}
		\dot{H} = -R I^2 \leq 0
	\end{equation}
	ist negativ semidefinit. Dass dies tatsächlich zur asymptotischen Stabilität reicht (und nicht nur zu Lya\-punov-Stabilität), folgt aus dem Satz von Barbalat oder der LaSalle-Invarianzprinzip, das besagt: Das System konvergiert zur größten invarianten Menge in $\{\dot{H} = 0\} = \{I = 0\}$. Diese ist der Ursprung. Diese mathematische Argumentation, die in der Kontrolltheorie alltäglich ist, hat ihren physikalischen Kern in dem in dieser Arbeit beobachteten Prinzip: Die Kopplung erzwingt die Nicht-Negativität der Energieform, und die Dissipation bewirkt die Konvergenz.

	Die Lyapunov-Theorie ermöglicht es, diese Erkenntnisse auf nichtlineare und zeitvariante Systeme zu verallgemeinern, für die keine geschlossene analytische Lösung existiert. Die Energiefunktion tritt dabei als Brücke zwischen der physikalischen Intuition und dem mathematischen Stabilitätsbeweis auf.

	% ---
	\subsubsection*{7. \quad Graphentheoretische Beschreibung elektrischer Netzwerke}

	Elektrische Netzwerke besitzen eine natürliche \textit{Graphenstruktur}: Die Knoten des Graphen entsprechen den Schaltungsknoten (Potenzialorten), die Kanten entsprechen den Bauelementen (Zweigen). Diese Beobachtung führt auf die \textit{algebraische Graphentheorie} der Netzwerke.

	Die Kirchhoffschen Gesetze haben in dieser Sprache eine elegante Matrizenformulierung. Sei $A$ die \textit{Inzidenzmatrix} des Graphen (Knoten $\times$ Zweige), dann lauten die Knotenpotentialität (KCL) und die Maschenspannung (KVL):
	\begin{equation}
		A \mathbf{i} = 0 \quad \text{(KCL)}, \qquad \mathbf{u} = A^T \boldsymbol{\varphi} \quad \text{(KVL)},
	\end{equation}
	wobei $\boldsymbol{\varphi}$ der Vektor der Knotenpotenziale ist. Die Gesamtleistung des Netzwerks ist:
	\begin{equation}
		P = \mathbf{u}^T \mathbf{i} = \boldsymbol{\varphi}^T A \mathbf{i} = 0,
	\end{equation}
	da $A \mathbf{i} = 0$. Dies ist der mathematische Beweis der Leistungserhaltung im Netzwerk -- ein Ergebnis, das in dieser Arbeit für den Transformator explizit gezeigt wurde.

	Der \textit{Laplace-Operator des Graphen} $\mathcal{L} = A \mathcal{Y} A^T$ (wobei $\mathcal{Y}$ die Diagonalmatrix der Admittanzen ist) ist positiv semidefinit -- eine direkte Konsequenz der Passivität der Bauelemente, die in dieser Arbeit an vielen Einzelbeispielen gezeigt wurde. Dieser Zusammenhang verbindet Netzwerktheorie, algebraische Graphentheorie und lineare Algebra auf tiefgründige Weise.

	Eine wichtige Erweiterung dieser graphentheoretischen Sichtweise ist die Beschreibung von neuronalen Netzen (sowohl biologischen als auch künstlichen) als gewichtete Graphen, in denen Signale analog zu Strömen und Spannungen fließen. Die in dieser Arbeit entwickelten Analogien legen nahe, dass Stabilitäts- und Energieprinzipien auch hier eine fundamentale Rolle spielen könnten.

	% ---
	\subsubsection*{8. \quad Diskrete Systeme und digitale Signalverarbeitung}

	Die in dieser Arbeit betrachteten Systeme sind ausnahmslos zeitkontinuierlich. In der modernen Signalverarbeitung und Regelungstechnik werden physikalische Systeme jedoch nahezu ausschließlich digital -- also zeitdiskret -- verarbeitet. Die Frage, ob und wie sich die hier entwickelten Prinzipien auf diskrete Systeme übertragen lassen, ist daher von großer praktischer Bedeutung.

	Der Ãœbergang vom kontinuierlichen zum diskreten System erfolgt durch die \textit{Bilineare Transformation} (Tustin-Methode) oder andere Diskretisierungsverfahren. Die zeitkontinuierliche Tiefpass-Ãœbertragungsfunktion
	\begin{equation}
		H(s) = \frac{1}{1 + \tau s}
	\end{equation}
	wird durch Substitution $s \to \frac{2}{T}\frac{z - 1}{z + 1}$ in die zeitdiskrete Ãœbertragungsfunktion
	\begin{equation}
		H(z) = \frac{T + 2\tau}{2(T + 2\tau)} \cdot \frac{1 + z^{-1}}{1 - \frac{2\tau - T}{2\tau + T} z^{-1}}
	\end{equation}
	umgewandelt. Die Komplementäritätseigenschaft von Tief- und Hochpassfilter ($|H_{TP}|^2 + |H_{HP}|^2 = 1$), die in Kapitel~15 und~16 für das kontinuierliche System gezeigt wurde, setzt sich im Diskreten fort -- eine Konsequenz der Unitarität der Systemmatrix, die wiederum eine Konsequenz der Energieerhaltung ist.

	In der digitalen Signalverarbeitung (DSP) werden diskrete Filter durch \textit{Differenzengleichungen} beschrieben, die die Stelle der Differentialgleichungen übernehmen. Für die numerische Simulation physikalischer Schaltungen stellt sich stets die Frage der \textit{numerischen Stabilität}: Ein Diskre\-tisierungsverfahren ist genau dann energieerhaltend (und damit stabil), wenn es die positiv definite Energieform des kontinuierlichen Systems erhält. Dies führt auf die Theorie der \textit{symmetrischen Runge-Kutta-Verfahren} und der \textit{geometrischen Numerik} -- ein aktives Forschungsgebiet, das Numerische Mathematik, Differential\-geometrie und Physik verbindet.

	% ---
	\subsubsection*{9. \quad Feldbeschreibung und partielle Differentialgleichungen}

	Die in dieser Arbeit betrachteten Systeme sind \textit{konzentrierte} (lumped) Parameter-Systeme: Jedes Bauelement wird durch einen einzigen skalaren Wert (Kapazität, Induktivität, Widerstand) charakterisiert. Diese Beschreibung ist eine Abstraktion der zugrundeliegenden physikalischen Realität, in der die Parameter räumlich verteilt sind.

	Der Übergang von konzentrierten zu \textit{verteilten} Parametern führt von gewöhnlichen Differentialgleichungen (ODEs) zu \textit{partiellen Differentialgleichungen} (PDEs). Das deutlichste Beispiel ist die Übertragungsleitung: Ein langes Koaxialkabel ist nicht durch einen einzelnen Widerstand, eine einzelne Kapazität und eine einzelne Induktivität beschreibbar, sondern durch die \textit{Telegraphengleichungen}:
	\begin{equation}
		\frac{\partial u}{\partial x} = -L'\frac{\partial i}{\partial t} - R'i, \qquad \frac{\partial i}{\partial x} = -C'\frac{\partial u}{\partial t} - G'u,
	\end{equation}
	wobei $L'$, $C'$, $R'$, $G'$ die Belags-Induktivität, -Kapazität, -Widerstand bzw.\ -Leitwert (pro Längeneinheit) sind. Diese Gleichungen sind die feldbeschreibende Verallgemeinerung des in dieser Arbeit betrachteten RLC-Seriengliedes.

	Auf der fundamentalsten Ebene werden elektromagnetische Phänomene durch die \textit{Maxwellschen Gleichungen} beschrieben:
	\begin{align}
		\nabla \times \mathbf{E} &= -\frac{\partial \mathbf{B}}{\partial t}, &
		\nabla \times \mathbf{H} &= \mathbf{J} + \frac{\partial \mathbf{D}}{\partial t}, \\
		\nabla \cdot \mathbf{D} &= \rho, &
		\nabla \cdot \mathbf{B} &= 0.
	\end{align}
	Die Energiedichte des elektromagnetischen Feldes ist:
	\begin{equation}
		w = \frac{1}{2}\left(\varepsilon_0 E^2 + \frac{B^2}{\mu_0}\right) \geq 0,
	\end{equation}
	eine direkte Verallgemeinerung der Energieausdrücke $\frac{1}{2}CU^2$ (elektrisches Feld im Kondensator) und $\frac{1}{2}LI^2$ (magnetisches Feld in der Spule). Das zentrale Prinzip dieser Arbeit -- die Nicht-Negativität der Energiedichte als Konsequenz der physikalischen Kopplungsrelationen -- ist also bereits in den Maxwellschen Gleichungen angelegt. Der Poynting-Satz, der die Energieerhaltung des elektromagnetischen Feldes beschreibt, ist das feldtheoretische Analogon der Energiebilanz des LC-Schwingkreises.

	% ---
	\subsubsection*{10. \quad Quantenmechanische Erweiterungen: Quanten-Schaltkreise}

	Ein faszinantes neues Forschungsgebiet ist die Ãœbertragung der in dieser Arbeit entwickelten Konzepte auf die Quantenmechanik -- insbesondere im Kontext der \textit{supraleitenden Quanten-Schaltkreise}, die die Basis moderner Quantencomputer bilden.

	In einem supraleitenden Josephson-Kontakt-Schaltkreis wird die klassische Energie des LC-Schwingkreises zum \textit{Hamilton-Operator}:
	\begin{equation}
		\hat{H} = \frac{(\hat{Q} - Q_g)^2}{2C} - E_J \cos\left(\frac{2\pi\hat{\Phi}}{\Phi_0}\right),
	\end{equation}
	wobei $\hat{Q}$ und $\hat{\Phi}$ konjugierte quantenmechanische Operatoren sind (analog zu Ort und Impuls), $E_J = \Phi_0 I_c / (2\pi)$ die Josephson-Energie und $\Phi_0 = h/(2e)$ das magnetische Flussquantum ist. Die Eigenwerte dieses Hamilton-Operators sind nicht-negativ (nach geeigneter Normierung), in direkter Analogie zur klassischen Energieform.

	Die Kopplung zwischen Ladung und Fluss -- $Q = C\,U$ und $\Phi = L\,I$ -- setzt sich in der Quantenmechanik als kanonische Kommutatorrelation fort:
	\begin{equation}
		[\hat{\Phi}, \hat{Q}] = i\hbar,
	\end{equation}
	analog zur Heisenbergschen Unschärferelation $[\hat{x}, \hat{p}] = i\hbar$ in der Mechanik. Das Thema dieser Arbeit -- die physikalische Kopplung zweier Größen als Fundament der mathematischen Struktur -- reicht also bis in die Quantenmechanik.

	Quanten-LC-Schwingkreise (Transmons, Fluxoniums etc.) sind die Bausteine aktueller Quantencomputer. Die Analyse und das Verständnis dieser Systeme stärkt wesentlich auf denselben konzeptionellen Grundlagen auf, die in dieser Arbeit für den klassischen LC-Schwingkreis entwickelt wurden: Energieerhaltung, Kopplungsrelationen und die mathematische Struktur der Differentialgleichungen (bzw.\ in der Quantenmechanik: der Schrödinger-Gleichung).

	% ---
	\subsubsection*{11. \quad Systemidentifikation und datengetriebene Modellierung}

	Ein weiterer bedeutsamer Ausblick gilt der Frage, wie die in dieser Arbeit entwickelten physikalischen Modelle in der Praxis aus Messdaten gewonnen -- oder überprüft -- werden können: die \textit{Systemidentifikation}.

	Gegeben sei eine unbekannte physikalische Schaltung, von der nur Eingangs- und Ausgangssignale bekannt sind. Das Ziel der Systemidentifikation ist es, ein mathematisches Modell $\hat{H}(j\omega)$ zu finden, das das Eingangs-Ausgangs-Verhalten der Schaltung möglichst genau beschreibt. Für die linearen, zeitinvarianten (LTI) Systeme dieser Arbeit kann dies durch \textit{Frequenzgang-Messung} und anschließende Modellanpassung geschehen.

	Das Wissen, dass alle in dieser Arbeit betrachteten passiven Systeme eine positiv reelle Übertragungsfunktion haben (d.\,h.\ $\mathrm{Re}(H(j\omega)) \geq 0$ für alle $\omega \geq 0$), ist dabei ein mächtiges Constraint für den Identifikationsprozess. Es verhindert, dass unvermeidliche Messfehler zu physikalisch unsinnigen Modellen führen, die Energie erzeugen würden. Die physikalische Konsistenz -- das Thema dieser Arbeit -- wird so zu einem Regularisierungsprinzip für die Datenanalyse.

	Im Zeitalter des maschinellen Lernens und der neuronalen Netze gewinnt dieser Gedanke neue Aktualität: \textit{Physics-Informed Neural Networks} (PINNs) sind neuronale Netze, die durch explizite Energiebedingungen regularisiert werden -- d.\,h.\ die physikalischen Kopplungsrelationen dieser Arbeit werden als Lernzwang in das Netzwerk eingebaut. Dies führt zu Modellen, die nicht nur die Daten gut fitten, sondern auch physikalisch konsistent sind.

	% ---
	\subsubsection*{12. \quad Philosophische Reflexion: Die Mathematisierbarkeit der Natur}

	Schließlich berührt das Thema dieser Arbeit eine fundamentale erkenntnistheoretische Frage, die Physiker und Mathematiker seit Jahrhunderten beschäftigt: \textit{Warum ist die Natur mathematisch beschreibbar?} Der Mathematiker Eugene Wigner sprach 1960 in einem berühmten Aufsatz von der "unerhörten Wirksamkeit der Mathematik in den Naturwissenschaften" (\textit{The unreasonable effectiveness of mathematics in the natural sciences}).

	Die in dieser Arbeit entwickelte Antwort -- oder zumindest ein Mosaikstein dazu -- ist folgende: Die Mathematisierbarkeit physikalischer Strukturen ist nicht zufällig. Sie folgt daraus, dass physikalische Gesetze \textit{Kopplungsrelationen} zwischen Zustandsgrößen darstellen, die die scheinbare Freiheit dieser Größen einschränken. Die Mathematik, die diese eingeschränkte Freiheit beschreibt, ist notwendigerweise reichhaltig: Sie enthält quadratische Formen, Differentialgleichungen, Matrizenstrukturen -- und all diese Strukturen spiegeln die physikalische Realität wider, weil sie aus ihr entstanden sind.

	Die Einsicht ist tiefgründig: Nicht die Physik passt sich der Mathematik an -- und nicht die Mathematik wird von der Physik erfunden. Vielmehr entfalten beide zusammen eine gemeinsame Struktur, die weder rein physikalisch noch rein mathematisch ist, sondern das Wesen beider vereint. Diese Struktur ist es, die die in dieser Arbeit immer wieder aufgetauchte "wechselseitige Ergänzung" von Physik und Mathematik erst möglich macht -- und die hoffen lässt, dass die in dieser Arbeit begonnene Suche nach gemeinsamen Mustern noch lange nicht abgeschlossen ist.

	% ---
	\subsubsection*{13. \quad Konkrete nächste Schritte und offene Forschungsfragen}

	Aus den vorangegangenen Überlegungen lassen sich konkrete weiterführende Untersuchungen ableiten, die an diese Arbeit anknüpfen könnten:

	\begin{enumerate}
		\item \textbf{Nichtlineare Erweiterung:} Systematische Untersuchung nichtlinearer Schaltungen (Transistorverstärker, Schaltnetzteile, chaotische Schaltungen nach Chua) auf die Gültigkeit der hier entwickelten Prinzipien. Insbesondere: Lässt sich stets eine Lyapunov-Funktion konstruieren, die der Energie des linearen Systems entspricht?

		\item \textbf{Mehrtor-Analyse:} Erweiterung der Übersichtstabelle (Tab.~\ref{tab:uebersicht}) auf Mehrtore. Hier ist die zentrale Frage: Welche Matrixeigenschaft tritt an die Stelle der skalaren Nicht-Negativität?

		\item \textbf{Diskrete Energieprinzipien:} Konstruktion energy-erhaltender Diskretisierungen für die hier betrachteten Systeme (Symplektische Integratoren für LC- und RLC-Schaltungen) und Vergleich ihrer numerischen Stabilität mit klassischen Verfahren.

		\item \textbf{Messtechnische Verifikation:} Experimenteller Aufbau der betrachteten Schaltungen und Messung der Übertragungsfunktionen, Energieverläufe und Zeitkonstanten. Vergleich mit den in dieser Arbeit analytisch hergeleiteten Ergebnissen.

		\item \textbf{Port-Hamiltonsche Formulierung:} Explizite Aufstellung der Port-Hamiltonschen Systemmatrizen $(J, R, G, H)$ für alle in Tab.~\ref{tab:uebersicht} aufgeführten Strukturen und Nachweis, dass die dort aufgeführten Eigenschaften (Nicht-Negativität, Komplementarität, Leistungserhaltung) direkte Konsequenzen der Port-Hamiltonschen Struktur sind.

		\item \textbf{Transistor und Verstärker als aktive Elemente:} Erweiterung der Betrachtung auf aktive Bauelemente (Transistoren, Operationsverstärker mit externer Versorgung). Hier gilt die Nicht-Negativität der Leistungsaufnahme nicht mehr global -- aber es gilt eine verallgemeinerte Energiebilanz unter Einbeziehung der Versorgungsquellen.

		\item \textbf{Stochastische Systeme und thermisches Rauschen:} Jeder reale Widerstand erzeugt thermisches Rauschen (Johnson-Nyquist-Rauschen). Die spektrale Leistungsdichte dieses Rauschens ist $S(f) = 4k_BT R$. Die Frage ist: Wie verändert das Hinzufügen dieses stochastischen Anteils die in dieser Arbeit entwickelten Konsistenzprinzipien? Dies führt auf die Theorie stochastischer Differentialgleichungen (Itô-Kalkül) im Kontext physikalischer Schaltungen.
	\end{enumerate}

	Diese offenen Fragen zeigen, dass die in dieser Arbeit begonnene Untersuchung der wechselseitigen Ergänzung von Physik und Mathematik bei weitem nicht abgeschlossen ist. Im Gegenteil: Die hier entwickelten Grundprinzipien -- Kopplungsrelation, Nicht-Negativität, Energiestruktur -- erweisen sich als äußerst fruchtbare Konzepte, die weit über die betrachteten Beispiele hinausweisen und Verbindungen zu den aktuell lebendigsten Forschungsgebieten der mathematischen Physik, der Kontrolltheorie und der Quantentechnologie herstellen.

	% ---
	\subsubsection*{14. \quad Berechenbarkeit und Beschreibbarkeit: Die Grenzen der physikalischen Modellierung}

	Die in Kapitel~\ref{sec:berechenbarkeit} entwickelten Überlegungen zur Berechenbarkeit und Beschreibbarkeit eröffnen eine fundamentale Perspektive, die weit über die betrachteten elektrischen Schaltungen hinausreicht. Es ist dies vielleicht der tiefgründigste Ausblick dieser Arbeit, weil er nicht eine Erweiterung des Rahmens beschreibt, sondern dessen Grenzen ausleuchtet.

	\textbf{Das Äquivalenzprinzip von Beschreibbarkeit und Berechenbarkeit.}
	Das zentrale Prinzip lautet: Was beschreibbar ist, ist berechenbar. Und: Wenn eine Funktion nicht berechenbar ist, dann ist sie auch nicht mehr beschreibbar -- denn alle beschreibbaren Funktionen sind berechenbar. Diese Äquivalenz ist kein Zirkelschluss, sondern eine strukturelle Aussage über die Natur mathematischer und physikalischer Erkenntnis. Eine endliche, regelgeleitete Beschreibung einer Funktion enthält stets implizit das Berechnungsverfahren. Wo keine Beschreibung existiert, existiert auch kein Verfahren.

	\textbf{Physikalische Phänomene jenseits der vollständigen Beschreibbarkeit.}
	Die in dieser Arbeit betrachteten physikalischen Strukturen sind ausnahmslos in wohldefinierten, stetigen und differenzierbaren Bereichen aus physikalischen Gründen beschreibbar. Dies ist kein Zufall: Die innere Ordnung dieser Systeme -- die Kopplungsrelationen $U = RI$, $Q = CU$, $\Phi = LI$ -- ist die physikalische Bedingung ihrer Beschreibbarkeit. Sie garantiert, dass die mathematische Beschreibung vollständig, widerspruchsfrei und berechenbar ist.

	Nicht alle physikalischen Phänomene besitzen diese Eigenschaft. Es gibt Naturerscheinungen, die nicht vollständig und allgemeingültig in einer geschlossenen mathematischen Funktion beschreibbar sind -- weil sie diskontinuierlich verlaufen, weil in ihnen Singularitäten auftreten, weil sie auf sensitive Anfangsbedingungen reagieren oder weil sie prinzipiell unvollständige Informationslagen erzeugen. Paradigmatische Beispiele sind:

	\begin{itemize}
		\item \textbf{Explosionen und Detonationen:} Hochgradig nichtlineare, diskontinuierliche Energiefreisetzung, die keine universell-geschlossene Beschreibungsfunktion erlaubt.
		\item \textbf{Vulkanausbrüche:} Geophysikalische Diskontinuitäten mit Sprüngen, Instabilitäten und Bifurkationen, die durch Gleichungen nur lokal und näherungsweise erfasst werden können.
		\item \textbf{Turbulente Strömungen:} Auch wenn die Navier-Stokes-Gleichungen die Physik korrekt formulieren, ist die Existenz regulärer Lösungen für turbulente Regime mathematisch nicht bewiesen -- die beschreibende Funktion ist nicht vollständig angebbar.
		\item \textbf{Deterministisches Chaos:} Systeme wie der Lorenz-Attraktor sind \textit{lokal} durch Differentialgleichungen beschreibbar, aber die Abbildung von Anfangszuständen auf zukünftige Zustände ist \textit{global} nicht geschlossen beschreibbar, weil beliebig kleine Unterschiede exponentiell verstärkt werden.
	\end{itemize}

	All diesen Phänomenen ist gemeinsam: Sie sind \textit{nicht immer kontinuierlich beschreibbar}. Ihr physikalischer Verlauf enthält Bereiche, in denen keine endliche, regelgeleitete Vorschrift existiert, die den weiteren Verlauf vollständig bestimmt.

	\textbf{Zustände, Teilfolgen und der letzte berechenbare Zustand.}
	Für den Fall, dass eine Funktion nicht mehr physikalisch oder mathematisch beschreibbar ist, aber berechenbar sein soll, gilt: Es ergibt sich dann eine Folge von Zuständen, für die dann am Ende eine Teilfolge von Zuständen existiert, dass wenn diese Teilfolge während der Berechnung betreten wird, dass es dann einen letzten Zustand in dieser Teilfolge gibt, ab dem die Maschine nicht mehr natürlich beschreibbar weiterrechnet.

	Dieser Sachverhalt beschreibt die strukturelle Situation an der Grenze der Beschreibbarkeit mit größtmöglicher Präzision. Eine Rechenmaschine durchläuft eine Folge von Zuständen -- solange die zugrundeliegende Funktion beschreibbar ist, schreitet sie regelgeleitet fort. Sobald der Bereich betreten wird, in dem das Phänomen nicht mehr beschreibbar ist, existiert innerhalb der Zustandsfolge eine Teilfolge, die in einen \textit{letzten Zustand} mündet: einen Zustand, ab dem kein natürlich beschreibbarer nächster Schritt mehr existiert. Die Berechnung kann nicht weiter nach einer durch die Physik vorgegebenen Beschreibungsregel fortgeführt werden.

	Dieses Bild ist eine strukturelle Aussage von grundlegender Bedeutung: Die Berechenbarkeit endet dort, wo die Beschreibbarkeit endet -- und die Beschreibbarkeit endet dort, wo das Naturphänomen aufhört, kontinuierlich und regelgeleitet zu sein.

	\textbf{Konsequenzen für das Programm dieser Arbeit.}
	Die hier entwickelte Perspektive schärft den Blick auf die Leistung und die Grenzen der in dieser Arbeit untersuchten Physik-Mathematik-Verbindung:

	\begin{enumerate}
		\item Die betrachteten elektronischen Strukturen sind \textit{ausgezeichnete} physikalische Systeme -- ausgezeichnet in dem Sinne, dass ihre Beschreibbarkeit, Berechenbarkeit und physikalische Konsistenz zusammenfallen. Dies ist kein Selbstverständlichkeit, sondern Folge ihrer inneren Ordnung.
		\item Die Grenzen dieser Ordnung sind nicht willkürlich, sondern durch die Physik selbst gegeben: Wo die stetige, differenzierbare, kopplungsgeregelte Dynamik endet, endet auch die vollständige Beschreibbarkeit -- und damit die vollständige Berechenbarkeit.
		\item Jede physikalische Theorie, die auf vollständiger mathematischer Beschreibbarkeit beruht, impliziert damit zugleich einen Gültigkeitsbereich: den Bereich, in dem das Phänomen tatsächlich durch die Beschreibungssprache erfasst wird. Jenseits dieses Bereichs -- bei der Explosion, beim Vulkanausbruch, in der voll entwickelten Turbulenz -- versagt nicht die Maschine. Es versagt die Beschreibung.
		\item Die in dieser Arbeit immer wieder formulierte wechselseitige Ergänzung von Physik und Mathematik ist deshalb nicht universell gültig, sondern an eine Bedingung geknüpft: die Beschreibbarkeit des Phänomens. Wo diese Bedingung erfüllt ist -- in den geordneten, kontinuierlichen, kopplungsgeregelten Strukturen dieser Arbeit -- entfaltet die Verbindung von Physik und Mathematik ihre volle Kraft. Wo sie nicht erfüllt ist, markiert sie eine Grenze der Erkenntnis.
	\end{enumerate}

	Der Ausblick auf die Grenzen der Beschreibbarkeit ist damit kein pessimistisches Schlusswort. Im Gegenteil: Er macht die Leistung der beschreibbaren Physik erst sichtbar. Die geordneten Systeme dieser Arbeit -- Schwingkreis, Filter, Transformator, Operationsverstärker -- sind Inseln der mathematischen Vollständigkeit in einem Meer physikalischer Phänomene, die sich der vollständigen Beschreibung entziehen. Sie zu verstehen, zu beschreiben und zu berechnen, ist die eigentliche Leistung der Verbindung von Physik und Mathematik.

	\begin{thebibliography}{99}
		
		\bibitem{shockley1949}
		W.~Shockley,
		\textit{The Theory of p-n Junctions in Semiconductors and p-n Junction Transistors},
		Bell System Technical Journal, vol.~28, no.~3, pp.~435--489, 1949.
		
		\bibitem{sze2006}
		S.~M.~Sze and K.~K.~Ng,
		\textit{Physics of Semiconductor Devices}, 3rd ed.
		Hoboken, NJ: Wiley-Interscience, 2006.
		
		\bibitem{streetman2015}
		B.~G.~Streetman and S.~K.~Banerjee,
		\textit{Solid State Electronic Devices}, 7th ed.
		Upper Saddle River, NJ: Pearson, 2015.
		
		\bibitem{neamen2011}
		D.~A.~Neamen,
		\textit{Semiconductor Physics and Devices: Basic Principles}, 4th ed.
		New York: McGraw-Hill, 2011.
		
		\bibitem{sedra2014}
		A.~S.~Sedra and K.~C.~Smith,
		\textit{Microelectronic Circuits}, 7th ed.
		Oxford: Oxford University Press, 2014.
		
		\bibitem{razavi2016}
		B.~Razavi,
		\textit{Design of Analog CMOS Integrated Circuits}, 2nd ed.
		New York: McGraw-Hill, 2016.
		
		\bibitem{pierret1996}
		R.~F.~Pierret,
		\textit{Semiconductor Device Fundamentals}.
		Reading, MA: Addison-Wesley, 1996.
		
		\bibitem{tsividis2011}
		Y.~Tsividis and C.~McAndrew,
		\textit{Operation and Modeling of the MOS Transistor}, 3rd ed.
		Oxford: Oxford University Press, 2011.
		
		\bibitem{weste2010}
		N.~H.~E.~Weste and D.~M.~Harris,
		\textit{CMOS VLSI Design: A Circuits and Systems Perspective}, 4th ed.
		Boston: Addison-Wesley, 2010.
		
		\bibitem{lundstrom2000}
		M.~Lundstrom,
		\textit{Fundamentals of Carrier Transport}, 2nd ed.
		Cambridge: Cambridge University Press, 2000.
		
		\bibitem{feynman2010}
		R.~P.~Feynman, R.~B.~Leighton, and M.~Sands,
		\textit{The Feynman Lectures on Physics}, vol.~II.
		New York: Basic Books, 2010.
		
		\bibitem{griffiths2017}
		D.~J.~Griffiths,
		\textit{Introduction to Electrodynamics}, 4th ed.
		Cambridge: Cambridge University Press, 2017.
		
		\bibitem{hayt2011}
		W.~H.~Hayt, J.~E.~Kemmerly, and S.~M.~Durbin,
		\textit{Engineering Circuit Analysis}, 8th ed.
		New York: McGraw-Hill, 2011.
		
		\bibitem{nilsson2014}
		J.~W.~Nilsson and S.~A.~Riedel,
		\textit{Electric Circuits}, 10th ed.
		Upper Saddle River, NJ: Pearson, 2014.
		
		\bibitem{boylestad2012}
		R.~L.~Boylestad and L.~Nashelsky,
		\textit{Electronic Devices and Circuit Theory}, 11th ed.
		Upper Saddle River, NJ: Pearson, 2012.
		
		\bibitem{horowitz1989}
		P.~Horowitz and W.~Hill,
		\textit{The Art of Electronics}, 3rd ed.
		Cambridge: Cambridge University Press, 2015.
		
		\bibitem{oppenheim2010}
		A.~V.~Oppenheim and A.~S.~Willsky,
		\textit{Signals and Systems}, 2nd ed.
		Upper Saddle River, NJ: Pearson, 1996.
		
		\bibitem{siebert1986}
		W.~McC.~Siebert,
		\textit{Circuits, Signals, and Systems}.
		Cambridge, MA: MIT Press, 1986.
		
		\bibitem{ogata2010}
		K.~Ogata,
		\textit{Modern Control Engineering}, 5th ed.
		Upper Saddle River, NJ: Prentice Hall, 2010.
		
		\bibitem{chua1969}
		L.~O.~Chua,
		\textit{Introduction to Nonlinear Network Theory}.
		New York: McGraw-Hill, 1969.
		
		\bibitem{gray2009}
		P.~R.~Gray, P.~J.~Hurst, S.~H.~Lewis, and R.~G.~Meyer,
		\textit{Analysis and Design of Analog Integrated Circuits}, 5th ed.
		Hoboken, NJ: Wiley, 2009.
		
		\bibitem{taur1998}
		Y.~Taur and T.~H.~Ning,
		\textit{Fundamentals of Modern VLSI Devices}, 2nd ed.
		Cambridge: Cambridge University Press, 2013.
		
		\bibitem{muller2002}
		R.~S.~Muller, T.~I.~Kamins, and M.~Chan,
		\textit{Device Electronics for Integrated Circuits}, 3rd ed.
		Hoboken, NJ: Wiley, 2002.
		
		\bibitem{kittel2004}
		C.~Kittel,
		\textit{Introduction to Solid State Physics}, 8th ed.
		Hoboken, NJ: Wiley, 2004.
		
		\bibitem{ashcroft1976}
		N.~W.~Ashcroft and N.~D.~Mermin,
		\textit{Solid State Physics}.
		New York: Holt, Rinehart and Winston, 1976.
		
		\bibitem{ioffe_database}
		A.~F.~Ioffe Institute,
		\textit{New Semiconductor Materials: Characteristics and Properties},
		Online database, St.~Petersburg, Russia, 2024.
		\url{http://www.ioffe.ru/SVA/NSM/Semicond/}
		
		\bibitem{ieee_std1149}
		IEEE Standard 1149.1,
		\textit{IEEE Standard Test Access Port and Boundary-Scan Architecture}.
		New York: IEEE, 2013.
		
		\bibitem{itrs2023}
		International Roadmap for Devices and Systems (IRDS),
		\textit{More Moore White Paper}.
		IEEE, 2023. \url{https://irds.ieee.org/}
		
		\bibitem{bergmann2014}
		L.~Bergmann and C.~Schaefer,
		\textit{Lehrbuch der Experimentalphysik, Band 6: Festk\"{o}rper}, 2.~Aufl.
		Berlin: De Gruyter, 2014.
		
		\bibitem{ibach2009}
		H.~Ibach and H.~L\"{u}th,
		\textit{Festk\"{o}rperphysik}, 8.~Aufl.
		Berlin: Springer, 2009.

		\bibitem{turing1936}
		A.~M.~Turing,
		\textit{On Computable Numbers, with an Application to the Entscheidungsproblem},
		Proceedings of the London Mathematical Society, vol.~42, pp.~230--265, 1936.

		\bibitem{church1936}
		A.~Church,
		\textit{An Unsolvable Problem of Elementary Number Theory},
		American Journal of Mathematics, vol.~58, no.~2, pp.~345--363, 1936.

		\bibitem{goedel1931}
		K.~G\"{o}del,
		\textit{\"{U}ber formal unentscheidbare S\"{a}tze der Principia Mathematica und verwandter Systeme~I},
		Monatshefte f\"{u}r Mathematik und Physik, vol.~38, pp.~173--198, 1931.

		\bibitem{rogers1967}
		H.~Rogers,
		\textit{Theory of Recursive Functions and Effective Computability}.
		Cambridge, MA: MIT Press, 1967.

		\bibitem{sipser2012}
		M.~Sipser,
		\textit{Introduction to the Theory of Computation}, 3rd ed.
		Boston: Cengage Learning, 2012.

		\bibitem{hopcroft2006}
		J.~E.~Hopcroft, R.~Motwani, and J.~D.~Ullman,
		\textit{Introduction to Automata Theory, Languages, and Computation}, 3rd ed.
		Upper Saddle River, NJ: Pearson, 2006.

		\bibitem{papadimitriou1994}
		C.~H.~Papadimitriou,
		\textit{Computational Complexity}.
		Reading, MA: Addison-Wesley, 1994.

		\bibitem{wigner1960}
		E.~P.~Wigner,
		\textit{The Unreasonable Effectiveness of Mathematics in the Natural Sciences},
		Communications on Pure and Applied Mathematics, vol.~13, no.~1, pp.~1--14, 1960.

		\bibitem{penrose1989}
		R.~Penrose,
		\textit{The Emperor's New Mind: Concerning Computers, Minds, and the Laws of Physics}.
		Oxford: Oxford University Press, 1989.

		\bibitem{penrose1994}
		R.~Penrose,
		\textit{Shadows of the Mind: A Search for the Missing Science of Consciousness}.
		Oxford: Oxford University Press, 1994.

		\bibitem{wolfram2002}
		S.~Wolfram,
		\textit{A New Kind of Science}.
		Champaign, IL: Wolfram Media, 2002.

		\bibitem{lorenz1963}
		E.~N.~Lorenz,
		\textit{Deterministic Nonperiodic Flow},
		Journal of the Atmospheric Sciences, vol.~20, no.~2, pp.~130--141, 1963.

		\bibitem{strogatz2014}
		S.~H.~Strogatz,
		\textit{Nonlinear Dynamics and Chaos: With Applications to Physics, Biology, Chemistry, and Engineering}, 2nd ed.
		Boulder, CO: Westview Press, 2014.

		\bibitem{gleick1987}
		J.~Gleick,
		\textit{Chaos: Making a New Science}.
		New York: Viking Penguin, 1987.

		\bibitem{mandelbrot1982}
		B.~B.~Mandelbrot,
		\textit{The Fractal Geometry of Nature}.
		New York: W.~H.~Freeman, 1982.

		\bibitem{kolmogorov1941}
		A.~N.~Kolmogorov,
		\textit{The Local Structure of Turbulence in Incompressible Viscous Fluid for Very Large Reynolds Numbers},
		Doklady Akademii Nauk SSSR, vol.~30, pp.~301--305, 1941.

		\bibitem{fefferman2006}
		C.~L.~Fefferman,
		\textit{Existence and Smoothness of the Navier-Stokes Equation},
		in: J.~Carlson, A.~Jaffe, A.~Wiles (Eds.), \textit{The Millennium Prize Problems}.
		Providence, RI: American Mathematical Society, 2006, pp.~57--67.

		\bibitem{temam1997}
		R.~Temam,
		\textit{Infinite-Dimensional Dynamical Systems in Mechanics and Physics}, 2nd ed.
		New York: Springer, 1997.

		\bibitem{arnol'd1989}
		V.~I.~Arnold,
		\textit{Mathematical Methods of Classical Mechanics}, 2nd ed.
		New York: Springer, 1989.

		\bibitem{ghrist2014}
		R.~Ghrist,
		\textit{Elementary Applied Topology}.
		CreateSpace, 2014.

		\bibitem{devaney2003}
		R.~L.~Devaney,
		\textit{An Introduction to Chaotic Dynamical Systems}, 2nd ed.
		Boulder, CO: Westview Press, 2003.

		\bibitem{smale1967}
		S.~Smale,
		\textit{Differentiable Dynamical Systems},
		Bulletin of the American Mathematical Society, vol.~73, no.~6, pp.~747--817, 1967.

		\bibitem{zeeman1977}
		E.~C.~Zeeman,
		\textit{Catastrophe Theory: Selected Papers 1972--1977}.
		Reading, MA: Addison-Wesley, 1977.

		\bibitem{thom1975}
		R.~Thom,
		\textit{Structural Stability and Morphogenesis: An Outline of a General Theory of Models}.
		Reading, MA: W.~A.~Benjamin, 1975.

		\bibitem{shannon1948}
		C.~E.~Shannon,
		\textit{A Mathematical Theory of Communication},
		Bell System Technical Journal, vol.~27, pp.~379--423 and 623--656, 1948.

		\bibitem{cover2006}
		T.~M.~Cover and J.~A.~Thomas,
		\textit{Elements of Information Theory}, 2nd ed.
		Hoboken, NJ: Wiley, 2006.

		\bibitem{chaitin1987}
		G.~J.~Chaitin,
		\textit{Algorithmic Information Theory}.
		Cambridge: Cambridge University Press, 1987.

		\bibitem{kolmogorov1965}
		A.~N.~Kolmogorov,
		\textit{Three Approaches to the Quantitative Definition of Information},
		Problems of Information Transmission, vol.~1, no.~1, pp.~1--7, 1965.

		\bibitem{born1927}
		M.~Born and R.~Oppenheimer,
		\textit{Zur Quantentheorie der Molekeln},
		Annalen der Physik, vol.~389, no.~20, pp.~457--484, 1927.

		\bibitem{heisenberg1927}
		W.~Heisenberg,
		\textit{\"{U}ber den anschaulichen Inhalt der quantentheoretischen Kinematik und Mechanik},
		Zeitschrift f\"{u}r Physik, vol.~43, pp.~172--198, 1927.

		\bibitem{dirac1930}
		P.~A.~M.~Dirac,
		\textit{The Principles of Quantum Mechanics}.
		Oxford: Clarendon Press, 1930.

		\bibitem{neumann1932}
		J.~von Neumann,
		\textit{Mathematische Grundlagen der Quantenmechanik}.
		Berlin: Springer, 1932.

		\bibitem{prigogine1984}
		I.~Prigogine and I.~Stengers,
		\textit{Order Out of Chaos: Man's New Dialogue with Nature}.
		New York: Bantam Books, 1984.

		\bibitem{nicolis1977}
		G.~Nicolis and I.~Prigogine,
		\textit{Self-Organization in Nonequilibrium Systems}.
		New York: Wiley, 1977.

		\bibitem{haken1983}
		H.~Haken,
		\textit{Synergetics: An Introduction}, 3rd ed.
		Berlin: Springer, 1983.

		\bibitem{mandelbrot1997}
		B.~B.~Mandelbrot,
		\textit{Fractals and Scaling in Finance: Discontinuity, Concentration, Risk}.
		New York: Springer, 1997.

		\bibitem{kantz1997}
		H.~Kantz and T.~Schreiber,
		\textit{Nonlinear Time Series Analysis}.
		Cambridge: Cambridge University Press, 1997.

		\bibitem{abramowitz1972}
		M.~Abramowitz and I.~A.~Stegun,
		\textit{Handbook of Mathematical Functions with Formulas, Graphs, and Mathematical Tables}.
		New York: Dover, 1972.

	\end{thebibliography}
	
\end{document}

